\documentclass{article}

% Language setting
% Replace `english' with e.g. `spanish' to change the document language
\usepackage[english]{babel}

% Set page size and margins
% Replace `letterpaper' with `a4paper' for UK/EU standard size

% Encoding
\usepackage[utf8]{inputenc}
\usepackage[T1]{fontenc}

% AMS
\usepackage{amsthm}
\usepackage{amsmath}
\usepackage{amsfonts}
\usepackage{amssymb}

% Graphics and tables
\usepackage{graphicx} 
\usepackage{graphbox}
\usepackage{float}
\usepackage{booktabs}
\usepackage{multirow}
\usepackage{rotating}
\usepackage{array}
\usepackage{xcolor}
\usepackage{subcaption}
\usepackage[ruled]{algorithm2e}

\usepackage{booktabs}
\usepackage{multirow}

% Comments
\usepackage{comment}

% New column types
\newcolumntype{C}[1]{>{\centering\arraybackslash}p{#1}}
\newcolumntype{R}[1]{>{\raggedleft\let\newline\\\arraybackslash\hspace{0pt}}m{#1}}
\newcolumntype{L}[1]{>{\raggedright\let\newline\\\arraybackslash\hspace{0pt}}m{#1}}

% Break citations
\usepackage{breakcites}

% No widows or orphans (use \nowidow[3] or \noclub[3] at the end of paragraph, respectively)
\usepackage{nowidow}

% For reference enumerations lists
\usepackage{enumitem}

% Margins setup
\usepackage[top=1.5cm,bottom=2cm,right=2.5cm,left=2.5cm]{geometry}
\usepackage{titling}

% References
\usepackage{natbib}

% Check the OS
\usepackage{ifplatform}

% Symbols
\usepackage{textcomp}
\usepackage{bbm}

% pdf version 1.7
\ifwindows
	\pdfoptionpdfminorversion 7
\fi

% No indentation in the beginning of the paragraph
% \setlength\parindent{0cm} 

% Allow math display page break
\allowdisplaybreaks

% Macros
\newcommand{\lp}{\left(}
\newcommand{\rp}{\right)}
\newcommand{\lc}{\left[}
\newcommand{\rc}{\right]}
\newcommand{\lb}{\left\{}
\newcommand{\rb}{\right\}}
\newcommand{\lf}{\left.}
\newcommand{\ri}{\right.}
%
\newcommand{\R}{\mathbb{R}}
\newcommand{\T}{\mathbb{T}}
\newcommand{\Sp}{\mathbb{S}}
\newcommand{\Z}{\mathbb{Z}}
\newcommand{\N}{\mathbb{N}}
%
\newcommand{\rd}{\mathrm{d}}
\newcommand{\cL}{\mathcal{L}}
\newcommand{\cA}{\mathcal{A}}
\newcommand{\bpsi}{\boldsymbol{\Psi}}
\newcommand{\buno}{\boldsymbol{1}}
\newcommand{\bl}{\boldsymbol{l}}
\newcommand{\bx}{\boldsymbol{x}}
\newcommand{\bz}{\boldsymbol{z}}
\newcommand{\bt}{\boldsymbol{t}}
\newcommand{\bs}{\boldsymbol{s}}
\newcommand{\be}{\boldsymbol{e}}
\newcommand{\bm}{\boldsymbol{m}}
\newcommand{\bn}{\boldsymbol{n}}
\newcommand{\by}{\boldsymbol{y}}
\newcommand{\bX}{\boldsymbol{X}}
\newcommand{\bY}{\boldsymbol{Y}}
\newcommand{\bZ}{\boldsymbol{Z}}
\newcommand{\bJ}{\boldsymbol{J}}
\newcommand{\bU}{\boldsymbol{U}}
\newcommand{\bR}{\boldsymbol{R}}
\newcommand{\bF}{\boldsymbol{F}}
\newcommand{\bmu}{\boldsymbol\mu}
\newcommand{\bu}{\boldsymbol{u}}
\newcommand{\bv}{\boldsymbol{v}}
\newcommand{\bb}{\boldsymbol{b}}
\newcommand{\bba}{\boldsymbol{a}}
\newcommand{\br}{\boldsymbol{r}}
\newcommand{\bh}{{\boldsymbol{h}}}
\newcommand{\bd}{\boldsymbol{d}}
\newcommand{\bk}{\boldsymbol{k}}
\newcommand{\bp}{\boldsymbol{p}}
\newcommand{\one}{\mathbf{1}}
\newcommand{\zero}{\mathbf{0}}
\newcommand{\two}{\mathbf{2}}
\newcommand{\bxi}{\boldsymbol\xi}
\newcommand{\bphi}{\boldsymbol\varphi}
\newcommand{\ba}{\boldsymbol\alpha}
\newcommand{\bbe}{\boldsymbol\beta}
\newcommand{\bbeta}{\boldsymbol\eta}
\newcommand{\bzeta}{\boldsymbol\zeta}
\newcommand{\bga}{\boldsymbol\gamma}
\newcommand{\bthe}{\boldsymbol\theta}
\newcommand{\btheta}{\boldsymbol\theta}
\newcommand{\bkappa}{\boldsymbol\kappa}
\newcommand{\bTheta}{\boldsymbol\Theta}
\newcommand{\Ical}{\mathcal{I}}
\newcommand{\bsigma}{\boldsymbol\sigma}
\newcommand{\bSigma}{\boldsymbol\Sigma}
\newcommand{\bLambda}{\boldsymbol\Lambda}
\newcommand{\bGamma}{\boldsymbol\Gamma}
\newcommand{\bDelta}{\boldsymbol\Delta}
\newcommand{\bO}{\mathbf{0}_q}
\newcommand{\bOd}{\mathbf{0}_{q-1}}
\newcommand{\Hcal}{\mathcal{H}}
\newcommand{\bHcal}{\boldsymbol{\mathcal{H}}}
\newcommand{\X}{\boldsymbol{\mathcal{X}}_{\bx,p}}
\newcommand{\W}{\boldsymbol{\mathcal{W}}_{\bx}}
\newcommand{\bnabla}{\boldsymbol\nabla}
\newcommand{\blambda}{\boldsymbol\lambda}
\newcommand{\bB}{\boldsymbol{B}}
\newcommand{\bK}{\boldsymbol{K}}
\newcommand{\bH}{\boldsymbol{H}}
\newcommand{\bA}{\boldsymbol{A}}
\newcommand{\bV}{\boldsymbol{V}}
\newcommand{\bI}{\boldsymbol{I}}
\newcommand{\bS}{\boldsymbol{S}}
\newcommand{\bP}{\boldsymbol{P}}
\newcommand{\bQ}{\boldsymbol{Q}}
\newcommand{\bT}{\boldsymbol{\mathcal{T}}}
\newcommand{\bM}{\boldsymbol{M}}
\newcommand{\bW}{\boldsymbol{W}}
\newcommand{\pr}{^{\prime}}
\newcommand{\cqfd}{\hfill $\square$}
%
\newcommand{\equald}{\stackrel{d}{=}}
\newcommand{\inlaw}{\stackrel{d}{\rightsquigarrow}}
\DeclareMathAlphabet\mathbfcal{OMS}{cmsy}{b}{n}
\DeclareMathOperator{\sech}{sech}
\DeclareMathOperator{\csch}{csch}
\DeclareMathOperator{\arcsec}{arcsec}
\DeclareMathOperator{\arccot}{arcCot}
\DeclareMathOperator{\arccsc}{arcCsc}
\DeclareMathOperator{\arccosh}{arcCosh}
\DeclareMathOperator{\arcsinh}{arcsinh}
\DeclareMathOperator{\arctanh}{arctanh}
\DeclareMathOperator{\arcsech}{arcsech}
\DeclareMathOperator{\arccsch}{arcCsch}
\DeclareMathOperator{\arccoth}{arcCoth} 
% More macros
\newcommand{\Xb}{\bX}
\newcommand{\Sb}{\mathbf{S}}
\newcommand{\thetab}{\boldsymbol\theta}
\newcommand{\etab}{\boldsymbol\eta}
\newcommand{\bnab}{\boldsymbol\nabla}
\newcommand{\xib}{\boldsymbol\xi}
\newcommand{\Vb}{\mathbf{V}}
\newcommand{\Qb}{\mathbf{Q}}
\newcommand{\Rb}{\mathbf{R}}
\newcommand{\Zb}{\mathbf{Z}}
\newcommand{\ab}{\mathbf{a}}
\newcommand{\ub}{\mathbf{u}}
\newcommand{\Tb}{\mathbf{T}}
\newcommand{\vb}{\mathbf{v}}
\newcommand{\xb}{\bx}
\newcommand{\yb}{\by}
\newcommand{\wb}{\mathbf{w}}
\newcommand{\Ab}{\mathbf{A}}
\newcommand{\Bb}{\bB}
\newcommand{\Hb}{\mathbf{H}}
\newcommand{\Ob}{\mathbf{O}}
\newcommand{\Cb}{\mathbf{C}}
\newcommand{\Db}{\mathbf{D}}
\newcommand{\Ub}{\mathbf{U}}
\newcommand{\Mb}{\mathbf{M}}
\newcommand{\Wb}{\mathbf{W}}
\newcommand{\Yb}{\mathbf{Y}}
\newcommand{\Lb}{\mathbf{L}}
\newcommand{\Gb}{\mathbf{G}}
\newcommand{\sigmad}{\sigma_{d}}
\newcommand{\sigmar}{\sigma_{\bd}}
\newcommand{\sigmai}{\sigma_{d_i}}
\newcommand{\sigmaim}{\sigma_{d_i-1}}
\newcommand{\sigmaj}{\sigma_{d_j}}
\newcommand{\sigmamj}{\sigma_{-d_j}}
\newcommand{\sigmamjk}{\sigma_{-(d_j,d_k)}}
\newcommand{\sigmak}{\sigma_{d_k}}
\newcommand{\sigmaell}{\sigma_{d_\ell}}
\newcommand{\sigmajm}{\sigma_{d_j-1}}
\newcommand{\sigmarm}{\sigma_{\bd-1}}
\newcommand{\sigmarmij}{\sigma_{d_1-1,\stackrel{i\neq j}{\ldots},d_r-1}}
\newcommand{\Sdr}{(\mathbb{S}^{d})^r}
\newcommand{\Sr}{\mathbb{S}^{\bd}}
\newcommand{\Srm}{\mathbb{S}^{\bd-1}}
\newcommand{\Sd}{\mathbb{S}^{d}}
\newcommand{\Sj}{\mathbb{S}^{d_j}}
\newcommand{\Suno}{\mathbb{S}^{d_1}}
\newcommand{\Sdos}{\mathbb{S}^{d_2}}
%\newcommand{\Smj}{\mathbb{S}^{-d_j}}
\newcommand{\Sk}{\mathbb{S}^{d_k}}
\newcommand{\Sell}{\mathbb{S}^{d_\ell}}
\newcommand{\Sjk}{\mathbb{S}^{d_j,d_k}}
\newcommand{\Smj}{\mathbb{S}^{\bd}_{-j}}
\newcommand{\Smjk}{\mathbb{S}^{\bd}_{-(j,k)}}
\newcommand{\Sjm}{\mathbb{S}^{d_j-1}}
\newcommand{\Sim}{\mathbb{S}^{d_i-1}}
\newcommand\independent{\protect\mathpalette{\protect\independenT}{\perp}}
\def\independenT#1#2{\mathrel{\rlap{$#1#2$}\mkern2mu{#1#2}}}

% Commands with arguments
\newcommand{\lrp}[1]{\left(#1\right)}
\newcommand{\lrc}[1]{\left[#1\right]}
\newcommand{\lrb}[1]{\left\{#1\right\}}
\newcommand{\lrpbig}[1]{\big(#1\big)}
\newcommand{\lrcbig}[1]{\big[#1\big]}
\newcommand{\lrbbig}[1]{\big\{#1\big\}}
\newcommand{\blrp}[1]{\big(#1\big)}
%
\newcommand{\Prob}[1]{\mathbb{P}\lb #1\rb}
%\newcommand{\E}[1]{\mathbb{E}\lp #1\rp}
\newcommand{\V}[1]{\mathbb{V}\mathrm{ar}\lc #1\rc}
\newcommand{\Es}[2]{\mathbb{E}_{#2}\lc #1\rc}
\newcommand{\Vs}[2]{\mathbb{V}\mathrm{ar}_{#2}\lc #1\rc}
\newcommand{\Cov}[2]{\mathbb{C}\mathrm{ov}\lc #1,#2\rc}
\newcommand{\supp}[1]{\mathrm{supp}\lp #1\rp}
\newcommand{\diag}[1]{\mathrm{diag}\lp #1\rp}
\newcommand{\cmod}[1]{\mathrm{cmod}\left(#1\right)}
%
\newcommand{\mse}[1]{\mathrm{MSE}\lrc{#1}}
\newcommand{\mise}[1]{\mathrm{MISE}\lrc{#1}}
\newcommand{\amise}[1]{\mathrm{AMISE}\lrc{#1}}
%
\newcommand{\df}[2]{\frac{\rd #1}{\rd #2}}
\newcommand{\dftwo}[2]{\frac{\rd^2 #1}{\rd #2^2}}
\newcommand{\pf}[2]{\frac{\partial #1}{\partial #2}}
\newcommand{\pftwo}[2]{\frac{\partial^2 #1}{\partial #2^2}}
\newcommand{\pfmix}[3]{\frac{\partial^2 #1}{\partial #2\partial #3}}
\newcommand{\norm}[1]{\left|\left| #1\right|\right|}
\newcommand{\abs}[1]{\left| #1\right|}
\newcommand{\tr}[1]{\mathrm{tr}\left[#1\right]}
\newcommand{\vect}[1]{\mathrm{vec}\left(#1\right)}
\newcommand{\vech}[1]{\mathrm{vech}\left(#1\right)}
\newcommand{\scalprod}[2]{\langle#1,#2\rangle}
%
\newcommand{\vlinel}[1]{\multicolumn{1}{|c}{#1}}
\newcommand{\vliner}[1]{\multicolumn{1}{c|}{#1}}
%
\newcommand{\alert}[1]{\textcolor{red}{#1}}

% Orders
\newcommand{\Om}[1]{\Omega_{#1}}
\newcommand{\om}[1]{\omega_{#1}}
\newcommand{\Iq}[2]{\int_{\Omega_{q}} #1\,\omega_{q}(d #2)}
\newcommand{\Iqr}[3]{\int_{\Omega_{q}\times\R} #1\,d #3\,\omega_{q}(d #2)}
\newcommand{\Iqp}[3]{\int_{\Omega_{q}\times\R^p} #1\,d #3\,\omega_{q}(d #2)}
\newcommand{\Iqm}[2]{\int_{\Omega_{q-1}} #1\,\omega_{q-1}(d #2)}
\newcommand{\Ir}[2]{\int_{\R} #1\,d #2}
\newcommand{\Iqq}[3]{\int_{\Omega_{q_1}\times\Omega_{q_2}} #1\,\omega_{q_2}(d #3)\,\omega_{q_1}(d #2)}
\newcommand{\Irr}[3]{\int_{\R\times\R} #1\,d #3\,d #2}
%
\DeclareFontFamily{OT1}{pzc}{}
\DeclareFontShape{OT1}{pzc}{m}{it}{<-> s * [1.10] pzcmi7t}{}
\DeclareMathAlphabet{\mathpzc}{OT1}{pzc}{m}{it}
\newcommand{\order}[1]{\mathpzc{o}\lp#1\rp}
\newcommand{\Order}[1]{\mathcal{O}\lp#1\rp}
\newcommand{\orderp}[1]{\mathpzc{o}_\mathbb{P}\lp#1\rp}
\newcommand{\Orderp}[1]{\mathcal{O}_\mathbb{P}\lp#1\rp}

% New counter
%\newcommand{\co}{\addtocounter{align}{1}\arabic{align}}

% New theorems
\newtheorem{definition}{Definition}[section]
\newtheorem{theorem}{Theorem}[section]
\newtheorem{corollary}{Corollary}[section]
\newtheorem{remark}{Remark}[section]
\newtheorem{proposition}{Proposition}[section]
\newtheorem{lemma}{Lemma}[section]
\newtheorem{algo}{Algorithm}[section]
\newtheorem{example}{Example}[section]
%\renewenvironment{proof}[1]{\textit{Proof#1.}}{\qed\\} 

% Roman enumerate
\renewcommand{\theenumi}{\textit{\roman{enumi}}}

% Percentage sign without typing \%
\newcommand{\pct}{\%}

% Definition
\newcommand{\defin}{:=}
\newcommand{\indef}{=:}


% % Useful packages
% \usepackage{amsmath}
% \usepackage{graphicx}
% \usepackage[colorlinks=true, allcolors=blue]{hyperref}
% \newtheorem{Definition}{Definition}[section] 
% \newtheorem{fig}{figure}[section] 
% \newtheorem{Theorem}{Theorem}[section] 
% \newtheorem{Proposition}{Proposition}[section] 
% \newtheorem{Corollary}{Corollary}[section] 
% \newtheorem{Example}{Example}[section] 
% \newtheorem{lemma}{Lemma}[section] 
% \newtheorem{Algorithm}{Algorithm}[section] 
% \newtheorem{Remark}{Remark}[section] 
% \newtheorem{Notation}{Notation}[section] 
% \newtheorem{Proof}{Proof}[section]

\title{On independence tests for polyspherical data and
their applications}
\author{Vinícius Litvinoff}

\begin{document}
\maketitle

\begin{abstract}
Your abstract.
\end{abstract}

\section{$I_{n1}$}

$I_{n1} = \int_{\Sr} [\hat{f}_{h_1, \dots, h_r}(\mathbf{x}) - \mathsf{E} (\hat{f}_{h_1, \dots, h_r}(\mathbf{x}))]^2 \mathrm{d} \mathbf{x}$.

\begin{align}
    I_{n1} &= \int_{\Sr} \Bigg[ \frac{1}{n} \sum_{i=1}^{n} \prod_{j=1}^{r} L_{\mathbf{h}} (x,X_j^{(i)}) - \mathsf{E} \Big( \frac{1}{n} \sum_{i=1}^{n} \prod_{j=1}^{r} L_{\mathbf{h}} (x,X_j^{(i)}) \Big) \Bigg]^2 \mathrm{d} \mathbf{x} \\
    &= \frac{1}{n^2} \int_{\Sr} \Bigg[ \sum_{i=1}^{n} \prod_{j=1}^{r} L_{\mathbf{h}} (x,X_j^{(i)}) - \mathsf{E} \Big( \sum_{i=1}^{n} \prod_{j=1}^{r} L_{\mathbf{h}} (x,X_j^{(i)}) \Big) \Bigg]^2 \mathrm{d} \mathbf{x} \\
    &= \frac{1}{n^2} \int_{\Sr} \Bigg[ \sum_{i=1}^{n} \prod_{j=1}^{r} L_{\mathbf{h}} (x,X_j^{(i)}) - \sum_{i=1}^{n} \mathsf{E} \Big( \prod_{j=1}^{r} L_{\mathbf{h}} (x,X_j^{(i)}) \Big) \Bigg]^2 \mathrm{d} \mathbf{x} \\
    &= \frac{1}{n^2} \int_{\Sr} \Bigg[ \sum_{i=1}^{n} \Bigg\{ \prod_{j=1}^{r} L_{\mathbf{h}} (x,X_j^{(i)}) - \mathsf{E} \Big( \prod_{j=1}^{r} L_{\mathbf{h}} (x,X_j^{(i)}) \Big) \Bigg\} \Bigg]^2 \mathrm{d} \mathbf{x} \\
    &= \frac{1}{n^2} \int_{\Sr} \Bigg[ \sum_{i=1}^{n} \Bigg\{ \prod_{j=1}^{r} L_{\mathbf{h}} (x,X_j^{(i)}) - \prod_{j=1}^{r} \mathsf{E} \Big( L_{\mathbf{h}} (x,X_j^{(i)}) \Big) \Bigg\} \Bigg]^2 \mathrm{d} \mathbf{x},
\end{align}
where the last equality comes from the null hypothesis of independence between the elements of the polysphere. Therefore, $I_{n1} = I_{n1}^{(1)} + I_{n1}^{(2)}$, where
\begin{align}
    I_{n1}^{(1)} &:= \frac{1}{n^2} \int_{\Sr} \Bigg[ \sum_{i=1}^{n} \Bigg\{ \prod_{j=1}^{r} L_{\mathbf{h}} (x,X_j^{(i)}) - \prod_{j=1}^{r} \mathsf{E} \Big( L_{\mathbf{h}} (x,X_j^{(i)}) \Big) \Bigg\}^2 \Bigg] \mathrm{d} \mathbf{x}; \\
    I_{n1}^{(2)} &:= \frac{2}{n^2} \int_{\Sr} \Bigg[ \sum_{1 \le i<k \le n} \Bigg\{ \prod_{j=1}^{r} L_{\mathbf{h}} (x,X_j^{(i)}) - \prod_{j=1}^{r} \mathsf{E} \Big( L_{\mathbf{h}} (x,X_j^{(i)}) \Big) \Bigg\} \\
    &\times \Bigg\{ \prod_{j=1}^{r} L_{\mathbf{h}} (x,X_j^{(k)}) - \prod_{j=1}^{r} \mathsf{E} \Big( L_{\mathbf{h}} (x,X_j^{(k)}) \Big) \Bigg\} \Bigg] \mathrm{d} \mathbf{x}.
\end{align}
The asymptotic distribution of $I_{n1} = I_{n1}^{(1)} + I_{n1}^{(2)}$ can be described through Lemma 1 in \cite{Garc_a_Portugu_es_2015}. In this case,
\begin{align}
    B_n^{-1/2} I_{n1} \xrightarrow{d} \mathcal{N}(0, 1),
\end{align}
where $B_n = n \mathsf{E} (\varphi_n^2(X_1)) + \frac{1}{2} n^2 \mathsf{E}(H_n^2(X_1,X_2))$ and $\varphi_n(X_1) = \frac{1}{n^2} \int_{\Sr} \Bigg[ \Bigg\{ \prod_{j=1}^{r} L_{\mathbf{h}} (x,X_j^{(1)}) - \prod_{j=1}^{r} \mathsf{E} \Big( L_{\mathbf{h}} (x,X_j^{(1)}) \Big) \Bigg\}^2 \Bigg] \mathrm{d} \mathbf{x}$.

\subsection{Proof using lemma 1}

Let $H_n(X_i,X_k) := \int_{\Sr} \Bigg[ \Bigg\{ \bar{L}_u (\bx,\bX_i) \Bigg\} \times \Bigg\{ \bar{L}_u (\bx,\bX_k) \Bigg\} \Bigg] \mathrm{d} \bx$.
\begin{align*}
    \mathsf{E} (H_n(X_1,X_2)) &= \mathsf{E} \Bigg( \int_{\Sr} \Bigg[ \Bigg\{ \bar{L}_u (\bx,\bX_1) \Bigg\} \times \Bigg\{ \bar{L}_u (\bx,\bX_2) \Bigg\} \Bigg] \mathrm{d} \bx \Bigg) \\
    &= \int_{\Sr} \mathsf{E} \Bigg( \Bigg[ \Bigg\{ \bar{L}_u (\bx,\bX_1) \Bigg\} \times \Bigg\{ \bar{L}_u (\bx,\bX_2) \Bigg\} \Bigg] \Bigg) \mathrm{d} \bx \\
    &= \int_{\Sr} \mathsf{E} \Bigg( \Bigg[ \Bigg\{ \bar{L}_u (\bx,\bX_1) \Bigg\} \Bigg) \mathsf{E} \Bigg( \Bigg\{ \bar{L}_u (\bx,\bX_2) \Bigg\} \Bigg] \Bigg) \mathrm{d} \bx.
\end{align*}

\begin{align*}
    \mathsf{E} (H_n^2(X_1,X_2)) &= \mathsf{E} \Bigg( \Bigg[  \int_{\Sr} \bar{L}_u (\bx,\bX_1) \bar{L}_u (\bx,\bX_2) \mathrm{d} \bx \Bigg]^2 \Bigg) \\
    &= \mathsf{E} \Bigg(  \int_{\Sr} \int_{\Sr} \Bigg\{ \bar{L}_u (\bx,\bX_1) \bar{L}_u (\bx,\bX_2) \Bigg\} \Bigg\{ \bar{L}_u (\by,\bX_1) \bar{L}_u (\by,\bX_2) \Bigg\}  \mathrm{d} \bx \mathrm{d} \by \Bigg) \\
    &= \int_{\Sr} \int_{\Sr} \mathsf{E} \Bigg( \Bigg\{ \bar{L}_u (\bx,\bX_1) \bar{L}_u (\bx,\bX_2) \Bigg\} \Bigg\{ \bar{L}_u (\by,\bX_1) \bar{L}_u (\by,\bX_2) \Bigg\} \Bigg) \mathrm{d} \bx \mathrm{d} \by \\
    &= \int_{\Sr} \int_{\Sr} \mathsf{E} \Bigg( \bar{L}_u (\bx,\bX_1) \bar{L}_u (\by,\bX_1) \Bigg) \mathsf{E} \Bigg( \bar{L}_u (\by,\bX_2) \bar{L}_u (\bx,\bX_2) \Bigg) \mathrm{d} \bx \mathrm{d} \by \\
    &= \int_{\Sr} \int_{\Sr} \mathsf{E}^2 \Bigg( \bar{L}_u (\bx,\bX_1) \bar{L}_u (\by,\bX_1) \Bigg) \mathrm{d} \bx \mathrm{d} \by.
\end{align*}

By a similar argument,
\begin{align*}
    \mathsf{E} (H_n^4(X_1,X_2)) &= \int_{\Sr} \dots \int_{\Sr} \mathsf{E}^2 \Bigg( \prod_{i=1}^{4} \bar{L}_u (\bx_i,\bX_1)  \Bigg) \prod_{i=1}^{4} \mathrm{d} \bx_i.
\end{align*}

$G_n(x,y) := \mathsf{E} (H_n(x,X_1),H_n(y,X_1))$. Therefore,
\begin{align*}
    G_n(s,t) &= \mathsf{E} \Bigg( \Bigg\{ \int_{\Sr} \bar{L}_u (\bx,\mathbf{s}) \bar{L}_u (\bx,\bX_1) \mathrm{d} \bx \Bigg\} \Bigg\{ \int_{\Sr} \bar{L}_u (\bx,\mathbf{t}) \bar{L}_u (\bx,\bX_1) \mathrm{d} \bx \Bigg\} \Bigg). \\
    \mathsf{E} (G_n^2(X_1,X_2)) &= \mathsf{E} \Bigg( \Bigg\{ \int_{\Sr} \bar{L}_u (\bx,X_1) \bar{L}_u (\bx,\bX_1) \mathrm{d} \bx \Bigg\}^2 \Bigg\{ \int_{\Sr} \bar{L}_u (\bx,X_2) \bar{L}_u (\bx,\bX_1) \mathrm{d} \bx \Bigg\}^2 \Bigg) \\
    &= \mathsf{E} \Bigg( \Bigg\{ \int_{\Sr} \int_{\Sr} \bar{L}_u (\bx_1,X_1) \bar{L}_u (\bx_1,\bX_1) \bar{L}_u (\bx_2,X_1) \bar{L}_u (\bx_2,\bX_1) \mathrm{d} \bx_1 \mathrm{d} \bx_2 \Bigg\} \\
    &\times \Bigg\{ \int_{\Sr} \int_{\Sr} \bar{L}_u (\bx_3,X_2) \bar{L}_u (\bx_3,\bX_1) \bar{L}_u (\bx_4,X_2) \bar{L}_u (\bx_4,\bX_1) \mathrm{d} \bx_3 \mathrm{d} \bx_4 \Bigg\} \Bigg) \\
    &= \mathsf{E} \Bigg( \Bigg\{ \int_{\Sr} \int_{\Sr} \bar{L}_u^2 (\bx_1,X_1) \bar{L}_u^2 (\bx_2,X_1) \mathrm{d} \bx_1 \mathrm{d} \bx_2 \Bigg\} \\
    &\times \Bigg\{ \int_{\Sr} \int_{\Sr} \bar{L}_u (\bx_3,X_2) \bar{L}_u (\bx_3,\bX_1) \bar{L}_u (\bx_4,X_2) \bar{L}_u (\bx_4,\bX_1) \mathrm{d} \bx_3 \mathrm{d} \bx_4 \Bigg\} \Bigg) \\
    &= \mathsf{E} \Bigg( \int_{\Sr} \int_{\Sr} \int_{\Sr} \int_{\Sr} \bar{L}_u^2 (\bx_1,X_1) \bar{L}_u^2 (\bx_2,X_1)  \\
    &\times  \bar{L}_u (\bx_3,X_2) \bar{L}_u (\bx_3,\bX_1) \bar{L}_u (\bx_4,X_2) \bar{L}_u (\bx_4,\bX_1) \mathrm{d} \bx_1 \mathrm{d} \bx_2 \mathrm{d} \bx_3 \mathrm{d} \bx_4 \Bigg) \\
    &= \int_{\Sr} \int_{\Sr} \int_{\Sr} \int_{\Sr} \mathsf{E} \Bigg( \bar{L}_u^2 (\bx_1,X_1) \bar{L}_u^2 (\bx_2,X_1)  \\
    &\times  \bar{L}_u (\bx_3,X_2) \bar{L}_u (\bx_3,\bX_1) \bar{L}_u (\bx_4,X_2) \bar{L}_u (\bx_4,\bX_1) \Bigg) \mathrm{d} \bx_1 \mathrm{d} \bx_2 \mathrm{d} \bx_3 \mathrm{d} \bx_4
\end{align*}

\begin{align*}
    \mathsf{E} (Z_{n1}^2) =&\; \mathsf{E} \Bigg( \Bigg(\int_{\Sr}  \bar{L}_{\bh}^2 (\bx,\bX_1)  \,\rd \bx \Bigg)^2 \Bigg) \\
    %&= \mathsf{E} \Bigg( \int_{\Sr} \int_{\Sr}  \bar{L}_{\bh}^2 (\bx,\bX_i) \bar{L}_{\bh}^2 (\by,\bX_i)  \,\rd \bx \,\rd \by \Bigg) \\
    %&= \int_{\Sr} \int_{\Sr} \mathsf{E} \Bigg(  \bar{L}_{\bh}^2 (\bx,\bX_i) \bar{L}_{\bh}^2 (\by,\bX_i) \Bigg)  \,\rd \bx \,\rd \by\\
    \leq &\; \mathsf{E} \Bigg( \Bigg(\int_{\Sr} 2\lrb{L^2_{\bh} (\bx,\bX_1) + \mathsf{E} [L_{\bh} (\bx,\bX_i)]^2 } \,\rd \bx \Bigg)^2 \Bigg) \\
    = &\; 4\mathsf{E} \Bigg( \Bigg(\int_{\Sr} L^2_{\bh} (\bx,\bX_1)\,\rd \bx  + \int_{\Sr}\mathsf{E} [L_{\bh} (\bx,\bX_i)]^2  \,\rd \bx \Bigg)^2 \Bigg) \\
    \leq &\; 8\mathsf{E} \Bigg( \Bigg(\int_{\Sr} L^2_{\bh} (\bx,\bX_1)\,\rd \bx\Bigg)^2  + \Bigg(\int_{\Sr}\mathsf{E} [L_{\bh} (\bx,\bX_i)]^2  \,\rd \bx \Bigg)^2 \Bigg) \\
    = &\; 8\mathsf{E} \Bigg( \Bigg(\int_{\Sr} L^2_{\bh} (\bx,\bX_1)\,\rd \bx\Bigg)^2 \Bigg)  + 8\mathsf{E} \Bigg( \Bigg(\int_{\Sr}\mathsf{E} [L_{\bh} (\bx,\bX_i)]^2  \,\rd \bx \Bigg)^2 \Bigg) \\
    = &\; 8\mathsf{E} \Bigg( \Bigg(\int_{\Sr} L^2_{\bh} (\bx,\bX_1)\,\rd \bx\Bigg)^2 \Bigg)  + 8 \Bigg(\int_{\Sr}\mathsf{E} [L_{\bh} (\bx,\bX_i)]^2  \,\rd \bx \Bigg)^2 \\
    = &\; 8\mathsf{E} ( O ( \rho ( \bh ) \lambda_d(L) )^2 ))  + 8 \Bigg(\int_{\Sr}\mathsf{E} [L_{\bh} (\bx,\bX_i)]^2  \,\rd \bx \Bigg)^2 \\
    = &\; 8 O (( \rho ( \bh ) \lambda_d(L) )^2 )  + 8 \Bigg(\int_{\Sr}\mathsf{E} [L_{\bh} (\bx,\bX_i)]^2  \,\rd \bx \Bigg)^2 \\
    = &\; 8 O (( \rho ( \bh ) \lambda_d(L) )^2 )  + 8 \Bigg(\int_{\Sr} [f(\boldsymbol{x})+h_n^2 b_d(L) \operatorname{tr}[\boldsymbol{\mathcal { T }} \bar{f}(\boldsymbol{x})]+o\left(h_n^2\right)]^2 \,\rd \bx \Bigg)^2 \\
    = &\; 8 O (( \rho ( \bh ) \lambda_d(L) )^2 )  + 8 \Bigg(\int_{\Sr} [(f(\boldsymbol{x})+h_n^2 b_d(L) \operatorname{tr}[\boldsymbol{\mathcal { T }} \bar{f}(\boldsymbol{x})])^2 +o\left(h_n^2\right)] \,\rd \bx \Bigg)^2 \\
    = &\; 8 O (( \rho ( \bh ) \lambda_d(L) )^2 )  + 8 \Bigg( o\left(h_n^2\right) + \int_{\Sr} [f(\boldsymbol{x})+h_n^2 b_d(L) \operatorname{tr}[\boldsymbol{\mathcal { T }} \bar{f}(\boldsymbol{x})]]^2 \,\rd \bx \Bigg)^2 \\
    = &\; 8 O (( \rho ( \bh ) \lambda_d(L) )^2 )  + 8 \Bigg\{ o\left(h_n^2\right) + \Bigg( \int_{\Sr} [f(\boldsymbol{x})+h_n^2 b_d(L) \operatorname{tr}[\boldsymbol{\mathcal { T }} \bar{f}(\boldsymbol{x})]]^2 \,\rd \bx \Bigg)^2 \Bigg\} \\
    = &\; O (( \rho ( \bh ) \lambda_d(L) )^2 ) + o\left(h_n^2\right)  + 8 \Bigg( \int_{\Sr} [f(\boldsymbol{x})+h_n^2 b_d(L) \operatorname{tr}[\boldsymbol{\mathcal { T }} \bar{f}(\boldsymbol{x})]]^2 \,\rd \bx \Bigg)^2.
    %= &\; 8\mathsf{E} \Bigg( \Bigg(\int_{\Sr} L^2_{\bh} (\bx,\bX_1)\,\rd \bx\Bigg)^2  + \Bigg(\int_{\Sr}[f(\boldsymbol{x})+h_n^2 b_d(L) \operatorname{tr}[\boldsymbol{\mathcal { T }} \bar{f}(\boldsymbol{x})]+o\left(h_n^2\right)]^2 \rd \bx \Bigg)
\end{align*}

Use that
\begin{align}
    \int_{\Sr} L_{\bh}^k(\bx,\by) g(\by)\,\rd\by\asymp \frac{c_{\bd,L}(\bh)^k}{c_{\bd,L^k}(\bh)} g(\bx) \asymp \frac{\lambda_{\bd}(L^k)\rho(\bh)}{[\lambda_{\bd}(L)\rho(\bh)]^k} g(\bx)=\frac{\lambda_{\bd}(L^k)}{\lambda_{\bd}(L)^k} \rho(\bh)^{-(k-1)}g(\bx). \label{eq:equivLh}
\end{align}
as $\bh\to\zero$, for $k\geq 1$ and an arbitrary function $g:\Sr\to\mathbb{R}$.

\subsection{Some further developments}

\begin{align*}
    H_n(\bX_1, \bX_2) &= \frac{2}{n^2} \int_{\Sr} \bar{L}_{\bh} (\bx,\bX_1) \bar{L}_{\bh} (\bx,\bX_2) \,\rd \bx.
\end{align*}

\begin{align*}
    H_n^4(\bX_1, \bX_2) &= \frac{16}{n^{16}} \int_{(\Sr)^4} \Bigg( \prod_{l=1}^{4} \bar{L}_{\bh} (\bx_l,\bX_1) \Bigg) \Bigg( \prod_{l=1}^{4} \bar{L}_{\bh} (\bx_l,\bX_2) \Bigg) \prod_{l=1}^{4} \,\rd 
    \bx_l.
\end{align*}

\begin{align*}
    \mathsf{E} (H_n^4(\bX_1, \bX_2)) &= \frac{16}{n^{16}} \int_{(\Sr)^4} \mathsf{E}^2 \Bigg( \prod_{l=1}^{4} \bar{L}_{\bh} (\bx_l,\bX_1) \Bigg) \prod_{l=1}^{4} \,\rd 
    \bx_l \\
    & \le \frac{16}{n^{16}} \int_{(\Sr)^4} \mathsf{E} \Bigg( \prod_{l=1}^{4} \bar{L}_{\bh}^2 (\bx_l,\bX_1) \Bigg) \prod_{l=1}^{4} \,\rd \bx_l \\
    &= \frac{16}{n^{16}} \mathsf{E} \Bigg[ \int_{(\Sr)^4} \Bigg( \prod_{l=1}^{4} \bar{L}_{\bh}^2 (\bx_l,\bX_1) \Bigg) \prod_{l=1}^{4} \,\rd \bx_l \Bigg] \\
    &= \frac{16}{n^{16}} \mathsf{E} \Bigg[ \Bigg( \int_{\Sr} \bar{L}_{\bh}^2 (\bx,\bX_1) ,\rd \bx \Bigg)^4 \Bigg] \\
    &= \frac{16}{n^{16}} \mathsf{E} (Z_{n1}^4).
\end{align*}

Applying cp inequality, we have that
\begin{align*}
    Z_{ni}^4 & \le \Bigg( \int_{\Sr} 2[L_{\bh}^2 (\bx,\bX_1) + \mathsf{E}^2 (L_{\bh} (\bx,\bX_1))] \rd \bx \Bigg)^4 \\
    &= 16 \Bigg( \int_{\Sr} L_{\bh}^2 (\bx,\bX_1) \rd \bx + \int_{\Sr} \mathsf{E}^2 (L_{\bh} (\bx,\bX_1)) \rd \bx \Bigg)^4 \\
    & \le 128 \Bigg( \Bigg( \int_{\Sr} L_{\bh}^2 (\bx,\bX_1) \rd \bx \Bigg)^4 + \Bigg( \int_{\Sr} \mathsf{E}^2 (L_{\bh} (\bx,\bX_1)) \rd \bx \Bigg)^4 \Bigg).
\end{align*}

Now applying Lemma XXX, 
\begin{align*}
    \int_{\Sr} L_{\bh}^2 (\bx,\bX_1) \rd \bx \asymp \frac{\lambda_{\bd}(L^2)}{\lambda_{\bd}(L)^2} \rho(\bh)^{-1}.
\end{align*}
\begin{align*}
    \int_{\Sr} \mathsf{E}^2 (L_{\bh} (\bx,\bX_i)) &= \int_{\Sr} \Bigg( \int_{\Sr} L_{\bh} (\bx,\bX_1)) f(\bx) \rd \bx \Bigg)^2 \rd \bx \\
    & \asymp \int_{\Sr} \Bigg( \frac{\lambda_{\bd}(L)}{\lambda_{\bd}(L)} f(\bx) \Bigg)^2 \rd \bx \\
    &= R(f).
\end{align*}

Therefore,
\begin{align*}
    \mathsf{E} (H_n^4(\bX_1, \bX_k)) &= \frac{16}{n^{16}} O(\rho(\bh)^{-4} + R(f)^4) \\
    &= O(n^{-16} \rho(\bh)^{-4} + n^{-16}).
\end{align*}

\subsubsection{Second term}

\begin{align*}
    H_n(x,X_1) &= \frac{2}{n^2} \int_{\Sr} \bar{L}_{\bh} (\bt,\bx) \bar{L}_{\bh} (\bt,\bX_1) \,\rd \bt.
\end{align*}

\begin{align*}
    G_n(x,y) &= \int_{\Sr} \lrb{\frac{2}{n^2} \int_{\Sr} \bar{L}_{\bh} (\bt,\bx) \bar{L}_{\bh} (\bt,\bx_1) \,\rd \bt} \lrb{\frac{2}{n^2} \int_{\Sr} \bar{L}_{\bh} (\bs,\by) \bar{L}_{\bh} (\bs,\bx_1) \,\rd \bs} f_{X_1}(x_1) \,\rd \bx_1 \\
    &= \frac{4}{n^4} \int_{\Sr} \int_{\Sr} \int_{\Sr} \bar{L}_{\bh} (\bt,\bx) \bar{L}_{\bh} (\bt,\bx_1) \bar{L}_{\bh} (\bs,\by) \bar{L}_{\bh} (\bs,\bx_1) f_{X_1}(x_1) \,\rd \bx_1 \,\rd \bt \,\rd \bs \\
    &= \frac{4}{n^4} \int_{\Sr} \int_{\Sr} \bar{L}_{\bh} (\bt,\bx) \bar{L}_{\bh} (\bs,\by) \mathsf{E} (\bar{L}_{\bh} (\bt,\bX_1) \bar{L}_{\bh} (\bs,\bX_1)) \,\rd \bt \,\rd \bs.
\end{align*}

Note that, by Cuachy-Schwarz inequality,
\begin{align*}
    \mathsf{E} (\bar{L}_{\bh} (\bt,\bX_1) \bar{L}_{\bh} (\bs,\bX_1)) & \le \sqrt{\mathsf{E}^2 (\bar{L}_{\bh} (\bt,\bX_1)) \mathsf{E}^2 (\bar{L}_{\bh} (\bs,\bX_1)))} \\
    &= \sqrt{\mathsf{Var} ({L}_{\bh} (\bt,\bX_1)) \mathsf{Var} ({L}_{\bh} (\bs,\bX_1)))}.
\end{align*}

Therefore,
\begin{align*}
    G_n(x,y) & \le \frac{4}{n^4} \int_{\Sr} \int_{\Sr} \bar{L}_{\bh} (\bt,\bx) \bar{L}_{\bh} (\bs,\by) \sqrt{\mathsf{Var} ({L}_{\bh} (\bt,\bX_1)) \mathsf{Var} ({L}_{\bh} (\bs,\bX_1)))} \,\rd \bt \,\rd \bs. \\
    G_n^2(x,y) & \le \frac{16}{n^{16}} \int_{(\Sr)^4} \bar{L}_{\bh} (\bt_1,\bx) \bar{L}_{\bh} (\bt_2,\by) \bar{L}_{\bh} (\bt_3,\bx) \bar{L}_{\bh} (\bt_4,\by) \sqrt{\prod_{i=1}^{4} \mathsf{Var} ({L}_{\bh} (\bt_i,\bX_1)) } \prod_{i=1}^{4} \,\rd \bt_i.
\end{align*}

Therefore,
\begin{align*}
    \mathsf{E} (G_n^2(X_1,X_2)) & \le \frac{16}{n^{16}} \int_{(\Sr)^4} \mathsf{E} \lrb{ \bar{L}_{\bh} (\bt_1,\bX_1) \bar{L}_{\bh} (\bt_2,\bX_2) \bar{L}_{\bh} (\bt_3,\bX_1) \bar{L}_{\bh} (\bt_4,\bX_2)} \sqrt{\prod_{i=1}^{4} \mathsf{Var} ({L}_{\bh} (\bt_i,\bX_1)) } \prod_{i=1}^{4} \,\rd \bt_i \\
    &= \frac{16}{n^{16}} \int_{(\Sr)^4} \mathsf{E} \lrb{ \bar{L}_{\bh} (\bt_1,\bX_1) \bar{L}_{\bh} (\bt_3,\bX_1) } \mathsf{E} \lrb{ \bar{L}_{\bh} (\bt_2,\bX_2) \bar{L}_{\bh} (\bt_4,\bX_2)} \sqrt{\prod_{i=1}^{4} \mathsf{Var} ({L}_{\bh} (\bt_i,\bX_1)) } \prod_{i=1}^{4} \,\rd \bt_i \\
    &= \frac{16}{n^{16}} \int_{(\Sr)^4} \mathsf{E} \lrb{ \bar{L}_{\bh} (\bt_1,\bX_1) \bar{L}_{\bh} (\bt_3,\bX_1) } \mathsf{E} \lrb{ \bar{L}_{\bh} (\bt_2,\bX_1) \bar{L}_{\bh} (\bt_4,\bX_1)} \sqrt{\prod_{i=1}^{4} \mathsf{Var} ({L}_{\bh} (\bt_i,\bX_1)) } \prod_{i=1}^{4} \,\rd \bt_i. \\
\end{align*}

Applying Cauchy-Schwartz inequality again, we obtain
\begin{align*}
    \mathsf{E} (G_n^2(X_1,X_2)) \le &\; \frac{16}{n^{16}} \int_{(\Sr)^4}  \sqrt{\mathsf{Var} ({L}_{\bh} (\bt_1,\bX_1)) \mathsf{Var} ({L}_{\bh} (\bt_3,\bX_1)))} \sqrt{\mathsf{Var} ({L}_{\bh} (\bt_2,\bX_1)) \mathsf{Var} ({L}_{\bh} (\bt_4,\bX_1)))} \\
    & \times \sqrt{\prod_{i=1}^{4} \mathsf{Var} ({L}_{\bh} (\bt_i,\bX_1)) } \prod_{i=1}^{4} \,\rd \bt_i \\
    =&\; \frac{16}{n^{16}} \int_{(\Sr)^4} \sqrt{\prod_{i=1}^{4} \mathsf{Var} ({L}_{\bh} (\bt_i,\bX_1)) } \sqrt{\prod_{i=1}^{4} \mathsf{Var} ({L}_{\bh} (\bt_i,\bX_1)) } \prod_{i=1}^{4} \,\rd \bt_i \\
    =&\; \frac{16}{n^{16}} \int_{(\Sr)^4} \prod_{i=1}^{4} \mathsf{Var} ({L}_{\bh} (\bt_i,\bX_1)) \prod_{i=1}^{4} \,\rd \bt_i \\
    =&\; \frac{16}{n^{16}} \int_{(\Sr)^4} \prod_{i=1}^{4} n \mathsf{Var} (\hat{f}(\bt_i;\bh)) \prod_{i=1}^{4} \,\rd \bt_i \\
    =&\; \frac{16}{n^{16}} \int_{(\Sr)^4} \prod_{i=1}^{4} n \lrb{\frac{v_{\bd}(L)}{n \rho(\bh)} f(\bt_i)(1 + o(1))} \prod_{i=1}^{4} \,\rd \bt_i \\
    =&\; \frac{16}{n^{16}} \int_{(\Sr)^4} \prod_{i=1}^{4} \frac{v_{\bd}(L)}{ \rho(\bh)} f(\bt_i)(1 + o(1)) \prod_{i=1}^{4} \,\rd \bt_i \\
    =&\; \frac{16}{n^{16}} \prod_{i=1}^{4} \int_{(\Sr)^4} \frac{v_{\bd}(L)}{ \rho(\bh)} f(\bt_i)(1 + o(1)) \,\rd \bt_i \\
    =&\; \frac{16}{n^{16}} \prod_{i=1}^{4} \frac{v_{\bd}(L)}{ \rho(\bh)} (1 + o(1)) \\
    =&\; \frac{16}{n^{16}} \frac{v_{\bd}^4(L)}{ \rho^4(\bh)} (1 + o(1)).
\end{align*}

As a consequence,
\begin{align*}
    \mathsf{E} (G_n^2(X_1,X_2)) &= O(n^{-16} \rho^{-4}(\bh) + n^{-16}).
\end{align*}

Defining $A_n \defin n^3 H_n^4(\bX_1, \bX_2) + n^4 \mathsf{E} (G_n^2(X_1,X_2))$, then, from the previous results,
\begin{align*}
    A_n &= O(n^{-12} \rho^{-4}(\bh) + n^{-12}) \\
    &= O(n^{-12} \rho^{-4}(\bh)).
\end{align*}

\subsubsection{Final ratio}

\begin{align*}
    B_n &= \frac{2}{n^2} \int_{\Sr} \int_{\Sr} \mathsf{E}^2 ( \bar{L}_{\bh} (\bx,\bX_1) \bar{L}_{\bh} (\by,\bX_1) ) \,\rd \bx \,\rd \by.
\end{align*}

Note that
\begin{align*}
    \mathsf{E} ( \bar{L}_{\bh} (\bx,\bX_1) \bar{L}_{\bh} (\by,\bX_1) ) &= \mathsf{E} ( L_{\bh} (\bx,\bX_1) L_{\bh} (\by,\bX_1) ) - \mathsf{E} (L_{\bh} (\bx,\bX_1)) \mathsf{E} (L_{\bh} (\by,\bX_1)).
\end{align*}

Furthermore,
\begin{align*}
    \mathsf{E} ( L_{\bh} (\bx,\bX_1) L_{\bh} (\by,\bX_1) ) &= \int_{\Sr} L_{\bh} (\bx,\bt) L_{\bh} (\by,\bt) f(\bt) \,\rd \bt \\
    & \asymp L_{\bh} (\bx,\by) f(\by).\\
    \mathsf{E} (L_{\bh} (\bx,\bX_1)) &= \int_{\Sr} L_{\bh} (\bx,\bt) f(\bt) \,\rd \bt \\
    & \asymp f(\bx).\\
    \mathsf{E} (L_{\bh} (\by,\bX_1)) &= \int_{\Sr} L_{\bh} (\by,\bt) f(\bt) \,\rd \bt \\
    & \asymp f(\by).\\
\end{align*}

As a consequence,
\begin{align*}
    \mathsf{E}^2 ( \bar{L}_{\bh} (\bx,\bX_1) \bar{L}_{\bh} (\by,\bX_1) ) & \asymp \lrb{L_{\bh} (\bx,\by) f(\by) - f(\bx)f(\by)}^2 \\
    &= L_{\bh}^2 (\bx,\by) f^2(\by) + f^2(\bx)f^2(\by) -2L_{\bh} (\bx,\by) f(\bx) f^2(\by).
\end{align*}

Let's integrate each term.
\begin{align*}
    \int_{\Sr} \int_{\Sr} L_{\bh}^2 (\bx,\by) f^2(\by) \,\rd \by \,\rd \bx & \asymp \int_{\Sr} \frac{\lambda_{\bd}(L^2)}{\lambda_{\bd}(L)^2} \rho(\bh)^{-1} f^2(\bx) \,\rd \bx \\
    &= \frac{\lambda_{\bd}(L^2)}{\lambda_{\bd}(L)^2} \rho(\bh)^{-1} R(f). \\
    \int_{\Sr} \int_{\Sr} f^2(\bx)f^2(\by) \,\rd \by \,\rd \bx &= R^2(f). \\
    \int_{\Sr} \int_{\Sr}-2L_{\bh} (\bx,\by) f(\bx) f^2(\by) \,\rd \by \,\rd \bx & \asymp -2 \int_{\Sr} f(\bx) f^2(\bx) \,\rd \bx\\
    &= -2 \int_{\Sr}f^3(\bx) \,\rd \bx.
\end{align*}

Finally,
\begin{align*}
    B_n \asymp \frac{2}{n^2} \lrb{\frac{\lambda_{\bd}(L^2)}{\lambda_{\bd}(L)^2} \rho(\bh)^{-1} R(f) + R^2(f) -2 \int_{\Sr}f^3(\bx) \,\rd \bx}.
\end{align*}

Notice that, therefore, $A_n = O(n^{-12} \rho^{-4}(\bh))$ and $B_n^2 \asymp n^{-4}(\rho(\bh)^{-1}+1)$, which entails that $A_n / B_n^2 \to 0$ as $n$ goes to infinity.

\section{$I_{n2}$}

\begin{align}
    \Big[ \prod_{i = 1}^{r} \hat{f}_{h_i}(\mathbf{x_i}) - \prod_{i = 1}^{r} \mathsf{E} (\hat{f}_{h_i} (\mathbf{x_i})) \Big]^2 &= \Bigg( \prod_{i = 1}^{r} \hat{f}_{h_i}(\mathbf{x_i}) \Bigg)^2 + \Bigg( \prod_{i = 1}^{r} \mathsf{E} (\hat{f}_{h_i} (\mathbf{x_i})) \Bigg)^2 - 2 \prod_{i = 1}^{r} \mathsf{E} (\hat{f}_{h_i} (\mathbf{x_i})) \prod_{i = 1}^{r} \hat{f}_{h_i}(\mathbf{x_i}) \\
    &= \Bigg( \prod_{i = 1}^{r} \hat{f}_{h_i}^2(\mathbf{x_i}) \Bigg) + \Bigg( \prod_{i = 1}^{r} \mathsf{E}^2 (\hat{f}_{h_i} (\mathbf{x_i})) \Bigg) - 2 \prod_{i = 1}^{r} \mathsf{E} (\hat{f}_{h_i} (\mathbf{x_i})) \hat{f}_{h_i}(\mathbf{x_i}).
\end{align}

\begin{align}
    \mathsf{E} \Bigg( \int \Big[ \prod_{i = 1}^{r} \hat{f}_{h_i}(\mathbf{x_i}) - \prod_{i = 1}^{r} \mathsf{E} (\hat{f}_{h_i} (\mathbf{x_i})) \Big]^2 \mathrm{d} \mathbf{x} \Bigg) &= \mathsf{E} \Bigg( \int  \Bigg( \prod_{i = 1}^{r} \hat{f}_{h_i}^2(\mathbf{x_i}) \Bigg) \mathrm{d} \mathbf{x} \Bigg)  + \mathsf{E} \Bigg( \int \Bigg( \prod_{i = 1}^{r} \mathsf{E}^2 (\hat{f}_{h_i} (\mathbf{x_i})) \Bigg) \mathrm{d} \mathbf{x} \Bigg) \\
    &- 2 \mathsf{E} \Bigg( \int \prod_{i = 1}^{r} \mathsf{E} (\hat{f}_{h_i} (\mathbf{x_i})) \hat{f}_{h_i} (\mathbf{x_i}) \mathrm{d} \mathbf{x} \Bigg) \\
    &= \mathsf{E} \Bigg( \prod_{i = 1}^{r} \int  \hat{f}_{h_i}^2(\mathbf{x_i}) \mathrm{d} \mathbf{x}_i \Bigg)  + \mathsf{E} \Bigg( \prod_{i = 1}^{r} \int \mathsf{E}^2 (\hat{f}_{h_i} (\mathbf{x_i})) \mathrm{d} \mathbf{x}_i \Bigg) \\
    &- 2 \mathsf{E} \Bigg( \prod_{i = 1}^{r} \int \mathsf{E} (\hat{f}_{h_i} (\mathbf{x_i})) \hat{f}_{h_i} (\mathbf{x_i}) \mathrm{d} \mathbf{x}_i \Bigg) \\
    &= \prod_{i = 1}^{r} \mathsf{E} \Bigg( \int  \hat{f}_{h_i}^2(\mathbf{x_i}) \mathrm{d} \mathbf{x}_i \Bigg)  + \prod_{i = 1}^{r} \mathsf{E} \Bigg( \int \mathsf{E}^2 (\hat{f}_{h_i} (\mathbf{x_i})) \mathrm{d} \mathbf{x}_i \Bigg) \\
     &- 2 \prod_{i = 1}^{r} \mathsf{E} \Bigg( \int \mathsf{E} (\hat{f}_{h_i} (\mathbf{x_i})) \hat{f}_{h_i} (\mathbf{x_i}) \mathrm{d} \mathbf{x}_i \Bigg) \\
     &= \prod_{i = 1}^{r} \int \mathsf{E} (  \hat{f}_{h_i}^2(\mathbf{x_i}) ) \mathrm{d} \mathbf{x}_i  + \prod_{i = 1}^{r} \int \mathsf{E} ( \mathsf{E}^2 (\hat{f}_{h_i} (\mathbf{x_i})) ) \mathrm{d} \mathbf{x}_i \\
     &- 2 \prod_{i = 1}^{r} \int \mathsf{E} ( \mathsf{E} (\hat{f}_{h_i} (\mathbf{x_i})) \hat{f}_{h_i} (\mathbf{x_i}) ) \mathrm{d} \mathbf{x}_i \\
     &= \prod_{i = 1}^{r} \int \mathsf{E} (  \hat{f}_{h_i}^2(\mathbf{x_i}) ) \mathrm{d} \mathbf{x}_i  + \prod_{i = 1}^{r} \int \mathsf{E}^2 (\hat{f}_{h_i} (\mathbf{x_i})) \mathrm{d} \mathbf{x}_i \\
     &- 2 \prod_{i = 1}^{r} \int \mathsf{E} (\hat{f}_{h_i} (\mathbf{x_i})) \mathsf{E} (\hat{f}_{h_i} (\mathbf{x_i}) ) \mathrm{d} \mathbf{x}_i \\
     &= \prod_{i = 1}^{r} \int \mathsf{E} (  \hat{f}_{h_i}^2(\mathbf{x_i}) ) \mathrm{d} \mathbf{x}_i  - \prod_{i = 1}^{r} \int \mathsf{E}^2 (\hat{f}_{h_i} (\mathbf{x_i})) \mathrm{d} \mathbf{x}_i.
\end{align}
Note that $\mathsf{E} (  \hat{f}_{h_i}^2(\mathbf{x_i}) ) = \mathsf{Var} (  \hat{f}_{h_i}(\mathbf{x_i}) ) + \mathsf{E}^2 (\hat{f}_{h_i} (\mathbf{x_i}))$. From results in \cite{Garc_a_Portugu_s_2013}, $\mathsf{Var} (  \hat{f}_{h_i}(\mathbf{x_i}) ) = O((nh_i^{q_i})^{-1})$ and $ \mathsf{E} (\hat{f}_{h_i} (\mathbf{x_i})) = f_{h_i} (\mathbf{x_i}) + O(h^2)$; as a consequence, $\mathsf{E}^2 (\hat{f}_{h_i} (\mathbf{x_i})) = f_{h_i}^2 (\mathbf{x_i}) + O(h^2)$, and, therefore, $\mathsf{E} (  \hat{f}_{h_i}^2(\mathbf{x_i}) ) = O((nh_i^{q_i})^{-1}) + f_{h_i}^2 (\mathbf{x_i}) + O(h^2)$.

Therefore,
\begin{align}
    \mathsf{E} \Bigg( \int \Big[ \prod_{i = 1}^{r} \hat{f}_{h_i}(\mathbf{x_i}) - \prod_{i = 1}^{r} \mathsf{E} (\hat{f}_{h_i} (\mathbf{x_i})) \Big]^2 \mathrm{d} \mathbf{x} \Bigg) &= \prod_{i = 1}^{r} \int [O((nh^{q_i})^{-1}) + f_i^2(x_i) + O(h^2)] \mathrm{d} \mathbf{x}_i  - \prod_{i = 1}^{r} \int [f_i^2(x_i) + O(h^2)] \mathrm{d} \mathbf{x}_i.
\end{align}
Now notice that
\begin{align}
    \prod_{i = 1}^{r} \int [O((nh^{q_i})^{-1}) + f_i^2(x_i) + O(h^2)] \mathrm{d} \mathbf{x}_i &= \prod_{i = 1}^{r} \Bigg\{ \int f_i^2(x_i) \mathrm{d} \mathbf{x}_i + O((nh^{q_i})^{-1}) + O(h^2) \Bigg\} \\
    &= \prod_{i = 1}^{r} \Bigg\{ \int f_i^2(x_i) \mathrm{d} \mathbf{x}_i \Bigg[1 + O((nh^{q_i})^{-1}) + O(h^2) \Bigg] \Bigg\} \\
    &= \Bigg\{ \prod_{i = 1}^{r} \int f_i^2(x_i) \mathrm{d} \mathbf{x}_i \Bigg\} \Bigg\{ \prod_{i = 1}^{r} [1 + O((nh^{q_i})^{-1}) + O(h^2)] \Bigg\} \\
    &= \Bigg\{ \prod_{i = 1}^{r} \int f_i^2(x_i) \mathrm{d} \mathbf{x}_i \Bigg\} [1 + O((nh^{q_i})^{-1}) + O(h^2)]^r \\
    &= \Bigg\{ \prod_{i = 1}^{r} \int f_i^2(x_i) \mathrm{d} \mathbf{x}_i \Bigg\} [1 + O(\mathsf{Max}((nh^{q_i})^{-1})+ h^2)].
\end{align}
By a similar argument,
\begin{align}
    \prod_{i = 1}^{r} \int [f_i^2(x_i) + O(h^2)] \mathrm{d} \mathbf{x}_i &= \Bigg\{ \prod_{i = 1}^{r} \int f_i^2(x_i) \mathrm{d} \mathbf{x}_i \Bigg\} [1 + O(h^2)].
\end{align}
Therefore,
\begin{align}
    \mathsf{E} \Bigg( \int \Big[ \prod_{i = 1}^{r} \hat{f}_{h_i}(\mathbf{x_i}) - \prod_{i = 1}^{r} \mathsf{E} (\hat{f}_{h_i} (\mathbf{x_i})) \Big]^2 \mathrm{d} \mathbf{x} \Bigg) &=\Bigg\{ \prod_{i = 1}^{r} \int f_i^2(x_i) \mathrm{d} \mathbf{x}_i \Bigg\} [1 + O(\mathsf{Max}((nh^{q_i})^{-1})+ h^2)] \\
    &- \Bigg\{ \prod_{i = 1}^{r} \int f_i^2(x_i) \mathrm{d} \mathbf{x}_i \Bigg\} [1 + O(h^2)] \\
    &= \Bigg\{ \prod_{i = 1}^{r} \int f_i^2(x_i) \mathrm{d} \mathbf{x}_i \Bigg\} [O(\mathsf{Max}((nh^{q_i})^{-1})+ h^2)].
\end{align}

\subsection{Proof using variance of $f_0$}

$I_{n2} := \int_{\Sr} \Big[ \prod_{i = 1}^{r} \hat{f}_{h_i}(\bx_i) - \prod_{i = 1}^{r} \mathsf{E} (\hat{f}_{h_i} (\bx_i)) \Big]^2 \mathrm{d} \bx$.

\begin{align*}
    \mathsf{E} (I_{n2}) &= \mathsf{E} \Bigg( \int_{\Sr} \Big[ \hat{f}_{\bh,0}(\bx) - \mathsf{E} (\hat{f}_{\bh,0}(\bx)) \Big]^2 \mathrm{d} \bx \Bigg) \\
    &= \int_{\Sr} \mathsf{E} \Bigg( \Big[ \hat{f}_{\bh,0}(\bx) - \mathsf{E} (\hat{f}_{\bh,0}(\bx)) \Big]^2 \Bigg) \mathrm{d} \bx \\
    &= \int_{\Sr} \mathsf{Var} (\hat{f}_{\bh,0}(\bx)) \mathrm{d} \bx \\
    &= \int_{\Sr} O((n \rho(\bh))^{-1}) \mathrm{d} \bx \\
    &= O((n \rho(\bh))^{-1}).
\end{align*}

\subsection{New proof}

Note that $\mathsf{E} (  \hat{f}_{h_i}^2(\mathbf{x_i}) ) = \mathsf{Var} (  \hat{f}_{h_i}(\mathbf{x_i}) ) + \mathsf{E}^2 (\hat{f}_{h_i} (\mathbf{x_i}))$. From results in \cite{Garc_a_Portugu_s_2013}, $\mathsf{Var} (  \hat{f}_{h_i}(\mathbf{x_i}) ) = O((nh_i^{q_i})^{-1})$ and $ \mathsf{E} (\hat{f}_{h_i} (\mathbf{x_i})) = f_{h_i} (\mathbf{x_i}) + O(h^2)$; as a consequence, $\mathsf{E}^2 (\hat{f}_{h_i} (\mathbf{x_i})) = f_{h_i}^2 (\mathbf{x_i}) + O(h^2)$, and, therefore, $\mathsf{E} (  \hat{f}_{h_i}^2(\mathbf{x_i}) ) = O((nh_i^{q_i})^{-1}) + f_{h_i}^2 (\mathbf{x_i}) + O(h^2)$.

\begin{align*}
        \mathsf{E} (I_{n2}) &= \mathsf{E} \Bigg( \int_{\Sr} \Big[ \hat{f}_{0}(\bx;\bh) - \mathsf{E} (\hat{f}_{0}(\bx;\bh)) \Big]^2 \,\rd \bx \Bigg) \\
        &= \int_{\Sr} \mathsf{E} \Bigg( \Big[ \hat{f}_{0}(\bx;\bh) - \mathsf{E} (\hat{f}_{0}(\bx;\bh)) \Big]^2 \Bigg) \,\rd \bx \\
        &= \int_{\Sr} \mathsf{Var} (\hat{f}_{0}(\bx;\bh)) \,\rd \bx \\
        &= \int_{\Sr} \mathsf{E} (\hat{f}_{0}^2(\bx;\bh)) - \mathsf{E}^2 (\hat{f}_{0}(\bx;\bh)) \,\rd \bx \\
        &= \int_{\Sr} \mathsf{E} (\prod_{j=1}^r \hat{f}_{j}^2(x^{(j)};h_j)) - \mathsf{E}^2 (\prod_{j=1}^r \hat{f}_{j}(x^{(j)};h_j)) \,\rd \bx \\
        &= \int_{\Sr} \Bigg\{ \prod_{j=1}^{r} [\mathsf{Var} (\hat{f}_j(x^{(j)};h_j)) + \mathsf{E}^2 (\hat{f}_j(x^{(j)};h_j))] - \prod_{j=1}^{r} \mathsf{E}^2 (\hat{f}_j(x^{(j)};h_j)) \Bigg\} \,\rd \bx \\
         &= \int_{\Sr} \Bigg\{ \prod_{j=1}^{r} [O((nh_j^{d_j})^{-1}) + f_{j}^2 (\bx; h_j) + O(h_j^2)] - \prod_{j=1}^{r} (f_{j}^2 (x^{(j)}; h_j) + O(h_j^2)) \Bigg\} \,\rd \bx \\
          &= \int_{\Sr} \Bigg\{ \Bigg[\prod_{j=1}^{r} f_{j}^2 (x^{(j)}; h_j) + O\Big(\sum_{j=1}^r [(nh_j^{d_j})^{-1} + h_j^2]\Big)\Bigg] - \Bigg[\prod_{j=1}^{r} f_{j}^2 (x^{(j)}; h_j) + O\Big(\sum_{j=1}^{r} h_j^2\Big)\Bigg] \Bigg\} \,\rd \bx \\
          &= O\Big(\sum_{j=1}^r [(nh_j^{d_j})^{-1} + h_j^2]\Big).
    \end{align*}

    $$\prod_{i=1}^n (a_i + b_i) = \sum_{\boldsymbol{k} \in \{0,1\}^n} \prod_{i=1}^n a_i^{1 -k_i} b_i^{k_i}$$

$$
\mathsf{E}^2 (\hat{f}_j(x^{(j)};h_j)=f_j(x^{(j)})+O(h_j^2)=O(1)
$$


\begin{align*}
    \prod_{j=1}^{r} [\mathsf{Var} (\hat{f}_j(x^{(j)};h_j)) + \mathsf{E}^2 (\hat{f}_j(x^{(j)};h_j))]=&\;\sum_{\boldsymbol{k} \in \{0,1\}^n} \prod_{j=1}^{r} \mathsf{Var} (\hat{f}_j(x^{(j)};h_j))^{1 -k_i} \mathsf{E}^2 (\hat{f}_j(x^{(j)};h_j))^{k_i}\\
    =&\;\prod_{j=1}^{r} \mathsf{E}^2 (\hat{f}_j(x^{(j)};h_j))\\
    &+\sum_{j=1}^r \mathsf{Var} (\hat{f}_j(x^{(j)};h_j)) \prod_{k=1,k\neq j}^{r} \mathsf{E}^2 (\hat{f}_k(x^{(k)};h_k))\\
    &+\cdots
\end{align*}

\begin{align*}
    \mathsf{E}[T_{n,2}]=\sum_{j=1}^r \frac{v_{d_j}(L)R(f_{0,-j})}{nh^{d_j}}+o\lrp{n^{-1}\sum_{j=1}^r h^{-d_j}},
\end{align*}
where $f_{0,-j}(\bx)=\frac{f_0(\bx)}{f_j(x^{(j)})}$.

$R(f_{0,-j})=\frac{R(f_0)}{R(f_j)}$ and $R(f)=\int f^2 $

\subsection{New proof (2)}

Let $R(f) \defin \int_{\Sr} f^2(\bx) \,\rd \bx$ and $f_{0,-j}(\bx) \defin \frac{f_0(\bx)}{f_j(\bx^{(j)})}$.

\begin{align*}
        \mathsf{E} (I_{n2}) &= \mathsf{E} \Bigg( \int_{\Sr} \Big[ \hat{f}_{0}(\bx;\bh) - \mathsf{E} (\hat{f}_{0}(\bx;\bh)) \Big]^2 \,\rd \bx \Bigg) \\
        &= \int_{\Sr} \mathsf{E} \Bigg( \Big[ \hat{f}_{0}(\bx;\bh) - \mathsf{E} (\hat{f}_{0}(\bx;\bh)) \Big]^2 \Bigg) \,\rd \bx \\
        &= \int_{\Sr} \mathsf{Var} (\hat{f}_{0}(\bx;\bh)) \,\rd \bx \\
        &= \int_{\Sr} \mathsf{E} (\hat{f}_{0}^2(\bx;\bh)) - \mathsf{E}^2 (\hat{f}_{0}(\bx;\bh)) \,\rd \bx \\
        &= \int_{\Sr} \mathsf{E} \Bigg(\prod_{j=1}^r \hat{f}_{j}^2(x^{(j)};h_j)\Bigg) - \mathsf{E}^2 \Bigg(\prod_{j=1}^r \hat{f}_{j}(x^{(j)};h_j)\Bigg) \,\rd \bx \\
        &= \int_{\Sr} \Bigg\{ \prod_{j=1}^{r} [\mathsf{Var} (\hat{f}_j(x^{(j)};h_j)) + \mathsf{E}^2 (\hat{f}_j(x^{(j)};h_j))] - \prod_{j=1}^{r} \mathsf{E}^2 (\hat{f}_j(x^{(j)};h_j)) \Bigg\} \,\rd \bx.
    \end{align*}

Note that, if any $b_j \to 0$ and $a_j = O(1)$, then
\begin{align*}
    \prod_{j=1}^r (a_j + b_j) &= \Bigg( \prod_{j=1}^{r} a_j\Bigg) + \sum_{j=1}^{r} b_j \Bigg( \prod_{k=1,k \ne j}^{r} a_k \Bigg) + o \Bigg( \sum_{j=1}^{r} b_j \Bigg\{ \prod_{k=1,k \ne j}^{r} a_k\Bigg\}\Bigg).
\end{align*}

Therefore,
\begin{align*}
    \prod_{j=1}^r (a_j + b_j) -  \prod_{j=1}^r a_j &= \sum_{j=1}^{r} b_j \Bigg( \prod_{k=1,k \ne j}^{r} a_k \Bigg) + o \Bigg( \sum_{j=1}^{r} b_j \Bigg\{ \prod_{k=1,k \ne j}^{r} a_k \Bigg\}\Bigg).
\end{align*}

Therefore, using the fact that $\mathsf{Var} (\hat{f}_j(x^{(j)};h_j)) = O((n h_j^{d_j})^{-1})$, and taking $a_j \defin \mathsf{E}^2 (\hat{f}_k(x^{(k)};h_k))$ and $b_j \defin \mathsf{Var} (\hat{f}_j(x^{(j)};h_j))$,
\begin{align*}
     \mathsf{E} (I_{n2}) =&\; \int_{\Sr} \Bigg\{ \Bigg[ \sum_{j=1}^{r} \mathsf{Var} (\hat{f}_j(x^{(j)};h_j)) \Bigg( \prod_{k=1,k \ne j}^{r} \mathsf{E}^2 (\hat{f}_k(x^{(k)};h_k)) \Bigg) \Bigg] \\
     & + o \Bigg( \sum_{j=1}^{r} \mathsf{Var} (\hat{f}_j(x^{(j)};h_j)) \Bigg\{ \prod_{k=1,k \ne j}^{r} \mathsf{E}^2 (\hat{f}_k(x^{(k)};h_k)) \Bigg\} \Bigg) \Bigg\} \,\rd \bx \\
     =&\; \int_{\Sr} \Bigg\{ \Bigg[ \sum_{j=1}^{r} \mathsf{Var} (\hat{f}_j(x^{(j)};h_j)) \Bigg( \prod_{k=1,k \ne j}^{r} \mathsf{E}^2 (\hat{f}_k(x^{(k)};h_k)) \Bigg) \Bigg] \\
     &+ o \Bigg( \sum_{j=1}^{r} (n h_j^{d_j})^{-1} f_j(x^{(j)}) v_{d_j}(L_j) \Bigg\{ \prod_{k=1,k \ne j}^{r} \mathsf{E}^2 (\hat{f}_k(x^{(k)};h_k)) \Bigg\}  \Bigg) \Bigg\} \,\rd \bx \\
     %&= \int_{\Sr} \Bigg\{ \Bigg[ \sum_{j=1}^{r} \mathsf{Var} (\hat{f}_j(x^{(j)};h_j)) \Bigg( \prod_{k=1,k \ne j}^{r} \mathsf{E}^2 (\hat{f}_k(x^{(k)};h_k)) \Bigg) \Bigg] + o \Bigg( n^{-1} \sum_{j=1}^{r} h_j^{-d_j} \Bigg) \Bigg\} \,\rd \bx \\
     =&\; \int_{\Sr} \Bigg\{ \sum_{j=1}^{r} \mathsf{Var} (\hat{f}_j(x^{(j)};h_j)) \Bigg( \prod_{k=1,k \ne j}^{r} \mathsf{E}^2 (\hat{f}_k(x^{(k)};h_k)) \Bigg) \Bigg\} \,\rd \bx + o \Bigg( n^{-1} \sum_{j=1}^{r} h_j^{-d_j} \Bigg) \\
     =&\; \sum_{j=1}^{r} \int_{\Sr} \Bigg\{ \mathsf{Var} (\hat{f}_j(x^{(j)};h_j)) \Bigg( \prod_{k=1,k \ne j}^{r} \mathsf{E}^2 (\hat{f}_k(x^{(k)};h_k)) \Bigg) \Bigg\} \,\rd \bx + o \Bigg( n^{-1} \sum_{j=1}^{r} h_j^{-d_j} \Bigg).
\end{align*}

Now notice that $\mathsf{E} (\hat{f}_k(x^{(k)};h_k)) = f_k(x^{(k)}) + O(h_j^2)$, therefore $\mathsf{E}^2 (\hat{f}_k(x^{(k)};h_k)) = f_k^2(x^{(k)}) + O(h_j^2)$. As a consequence,
\begin{align*}
    \prod_{k=1,k \ne j}^{r} \mathsf{E}^2 (\hat{f}_k(x^{(k)};h_k)) = \prod_{k=1,k \ne j}^{r} f_k^2(x^{(k)}) + O \Bigg( \sum_{j=1}^{r} h_j^2 \Bigg)= f_{0,-j}(\bx)^2+ O \Bigg( \sum_{j=1}^{r} h_j^2 \Bigg).
\end{align*}

Furthermore,
\begin{align*}
    \mathsf{Var} (\hat{f}_j(x^{(j)};h_j)) &= \frac{v_{d_j}(L)}{n h_j^{d_j}} f_j(x^{(j)}) (1 + o(1)).
\end{align*}

Therefore,
\begin{align*}
    \mathsf{E} (I_{n2}) &= \sum_{j=1}^{r} \int_{\Sr} \Bigg\{ \Bigg[ \frac{v_{d_j}(L)}{n h_j^{d_j}} f_j(x^{(j)})(1 + o(1)) \Bigg] \Bigg[ f_{0,-j}(\bx)^2+ O \Bigg( \sum_{j=1}^{r} h_j^2 \Bigg) \Bigg] \Bigg\} \,\rd \bx + o \Bigg( n^{-1} \sum_{j=1}^{r} h_j^{-d_j} \Bigg)\\
    &=: \sum_{j=1}^{r} \int_{\Sr} \theta_j \,\rd \bx + o \Bigg( n^{-1} \sum_{j=1}^{r} h_j^{-d_j} \Bigg).
\end{align*}
Then,
\begin{align*}
    \theta_j =&\; \frac{v_{d_j}(L)}{n h_j^{d_j}} f_j(x^{(j)}) f_{0,-j}(\bx)^2\\
    &+ \frac{v_{d_j}(L)}{n h_j^{d_j}} f_j(x^{(j)}) O \Bigg( \sum_{j=1}^{r} h_j^2 \Bigg) \\
    &+ o((n h_j^{d_j})^{-1}) f_{0,-j}(\bx)^2\\
    &+ o((n h_j^{d_j})^{-1}) O \Bigg( \sum_{j=1}^{r} h_j^2 \Bigg).
\end{align*}
Note that the last term of $\theta_j$ is negligible. Further,
\begin{align*}
    \int_{\Sr} \theta_j \,\rd \bx =&\; \int_{\Sr} \frac{v_{d_j}(L)}{n h_j^{d_j}} f_j(x^{(j)}) f_{0,-j}(\bx)^2 \,\rd \bx \\
    &+ \frac{v_{d_j}(L)}{n h_j^{d_j}} O \Bigg( \sum_{j=1}^{r} h_j^2 \Bigg) \\
    &+ o((n h_j^{d_j})^{-1})\\
    =&\; \frac{v_{d_j}(L)}{n h_j^{d_j}} \int_{\mathbb{S}^{d_j}} f_j(x^{(j)})\,\rd \bx^{(j)}  \int_{\mathbb{S}^{-d_j}} f_{0,-j}(\bx)^2 \,\rd \bx^{(-j)} \\
    &+ o((n h_j^{d_j})^{-1})\\
    =&\; \frac{v_{d_j}(L)}{n h_j^{d_j}}  R(f_{0,-j})+ o((n h_j^{d_j})^{-1}).
\end{align*}

Therefore,
\begin{align*}
    \mathsf{E} (I_{n2}) &= \sum_{j=1}^{r} \left\{\frac{v_{d_j}(L)}{n h_j^{d_j}}  R(f_{0,-j})+ o((n h_j^{d_j})^{-1})\right\}+ o \Bigg( n^{-1} \sum_{j=1}^{r} h_j^{-d_j} \Bigg)\\
    &=\sum_{j=1}^{r} \frac{v_{d_j}(L)}{n h_j^{d_j}}  R(f_{0,-j})+ o \Bigg( n^{-1} \sum_{j=1}^{r} h_j^{-d_j} \Bigg).
\end{align*}

\section{Variance of $T_{n,2}$}

Recall that
\begin{align*}
    T_{n,2} &= \int_{\Sr} [ \hat{f}_{0}(\bx;\bh) - \mathsf{E} (\hat{f}_{0}(\bx;\bh)) ]^2 \,\rd \bx.
\end{align*}

Therefore,
\begin{align*}
    \mathsf{E} (T_{n,2}^2) &= \mathsf{E} \lrb{ \lrb{ \int_{\Sr} [ \hat{f}_{0}(\bx;\bh) - \mathsf{E} (\hat{f}_{0}(\bx;\bh)) ]^2 \,\rd \bx }^2 } \\
    &= \mathsf{E} \lrb{ \lrb{ \int_{\Sr} [ \hat{f}_{0}(\bx;\bh) - \mathsf{E} (\hat{f}_{0}(\bx;\bh)) ]^2 \,\rd \bx } \lrb{ \int_{\Sr} [ \hat{f}_{0}(\by;\bh) - \mathsf{E} (\hat{f}_{0}(\by;\bh)) ]^2 \,\rd \by } } \\
    &= \mathsf{E} \lrb{ \int_{\Sr} \int_{\Sr} [ \hat{f}_{0}(\bx;\bh) - \mathsf{E} (\hat{f}_{0}(\bx;\bh)) ]^2 [ \hat{f}_{0}(\by;\bh) - \mathsf{E} (\hat{f}_{0}(\by;\bh)) ]^2 \,\rd \bx \,\rd \by } \\
    &= \int_{\Sr} \int_{\Sr} \mathsf{E} \lrb{ [ \hat{f}_{0}(\bx;\bh) - \mathsf{E} (\hat{f}_{0}(\bx;\bh)) ]^2 [ \hat{f}_{0}(\by;\bh) - \mathsf{E} (\hat{f}_{0}(\by;\bh)) ]^2 } \,\rd \bx \,\rd \by.
\end{align*}

Applying cp inequality,
\begin{align*}
    \mathsf{E} (T_{n,2}^2) & \le \int_{\Sr} \int_{\Sr} \mathsf{E} \lrb{ 4[ \hat{f}_{0}^2(\bx;\bh) + \mathsf{E}^2 (\hat{f}_{0}(\bx;\bh)) ] [ \hat{f}_{0}^2(\by;\bh) + \mathsf{E}^2 (\hat{f}_{0}(\by;\bh)) ] } \,\rd \bx \,\rd \by \\
    &= 4 \int_{\Sr} \int_{\Sr} \mathsf{E} \lrb{ \hat{f}_{0}^2(\bx;\bh)\hat{f}_{0}^2(\by;\bh) + \hat{f}_{0}^2(\bx;\bh)\mathsf{E}^2 (\hat{f}_{0}(\by;\bh)) + \mathsf{E}^2 (\hat{f}_{0}(\bx;\bh))\hat{f}_{0}^2(\by;\bh) + \mathsf{E}^2 (\hat{f}_{0}(\bx;\bh))\mathsf{E}^2 (\hat{f}_{0}(\by;\bh)) } \,\rd \bx \,\rd \by. \\
\end{align*}

We need to compute $\mathsf{E} (\hat{f}_{0}^2(\bx;\bh)\hat{f}_{0}^2(\by;\bh))$. Note that
\begin{align*}
    \hat{f}_{0}(\bx;\bh)\hat{f}_{0}(\by;\bh) &= \lrb{\frac{1}{n^r} \prod_{j=1}^{r} \sum_{i=1}^{n} L_{h_j} (\bx^{(j)}, \bX_i^{(j)})} \lrb{\frac{1}{n^r} \prod_{l=1}^{r} \sum_{k=1}^{n} L_{h_l} (\by^{(l)}, \bX_k^{(l)})} \\
    &= \frac{1}{n^{2r}} \prod_{j=1}^{r} \lrb{ \lrb{ \sum_{i=1}^{n} L_{h_j} (\bx^{(j)}, \bX_i^{(j)})} \lrb{ \sum_{k=1}^{n} L_{h_l} (\by^{(j)}, \bX_k^{(j)})} }.
\end{align*}

Therefore,
\begin{align*}
    \hat{f}_{0}^2(\bx;\bh)\hat{f}_{0}^2(\by;\bh) &= \frac{1}{n^{4r}} \prod_{j=1}^{r} \sum_{i_1,\dots,i_4 = 1}^{n} L_{h_j} (\bx^{(j)}, \bX_{i_1}^{(j)}) L_{h_j} (\bx^{(j)}, \bX_{i_1}^{(j)})
    L_{h_j} (\by^{(j)}, \bX_{i_3}^{(j)}) L_{h_j} (\by^{(j)}, \bX_{i_4}^{(j)}).
\end{align*}

\begin{align*}
    \mathsf{E} (\hat{f}_{0}^2(\bx;\bh)\hat{f}_{0}^2(\by;\bh)) &= \frac{1}{n^{4r}} \prod_{j=1}^{r} \sum_{i_1,\dots,i_4 = 1}^{n} \mathsf{E} \lrb{ L_{h_j} (\bx^{(j)}, \bX_{i_1}^{(j)}) L_{h_j} (\bx^{(j)}, \bX_{i_2}^{(j)})
    L_{h_j} (\by^{(j)}, \bX_{i_3}^{(j)}) L_{h_j} (\by^{(j)}, \bX_{i_4}^{(j)}) } \\
    &= \frac{1}{n^{3r}} \prod_{j=1}^{r} \sum_{i_2,\dots,i_4 = 1}^{n} \mathsf{E} \lrb{ L_{h_j} (\bx^{(j)}, \bX_{1}^{(j)}) L_{h_j} (\bx^{(j)}, \bX_{i_2}^{(j)})
    L_{h_j} (\by^{(j)}, \bX_{i_3}^{(j)}) L_{h_j} (\by^{(j)}, \bX_{i_4}^{(j)}) }
\end{align*}

Note that, when $i_2 = i_3 = i_4 = 1$,
\begin{align*}
     \mathsf{E} \lrb{ L_{h_j} (\bx^{(j)}, \bX_{1}^{(j)}) L_{h_j} (\bx^{(j)}, \bX_{i_2}^{(j)})
    L_{h_j} (\by^{(j)}, \bX_{i_3}^{(j)}) L_{h_j} (\by^{(j)}, \bX_{i_4}^{(j)}) } &= \mathsf{E} \lrb{ L_{h_j}^2 (\bx^{(j)}, \bX_{1}^{(j)})
    L_{h_j}^2 (\by^{(j)}, \bX_{1}^{(j)}) }
\end{align*}

Using Lemma B4, we can deduce that
\begin{align*}
    \mathsf{E} \lrb{ L_{h_j}^2 (\bx^{(j)}, \bX_{1}^{(j)})
    L_{h_j}^2 (\by^{(j)}, \bX_{1}^{(j)}) } &= \int_{\Sr}  L_{h_j}^2 (\bx^{(j)}, \bt)
    L_{h_j}^2 (\by^{(j)}, \bt) f_j(\bt) \,\rd \bt \\
    & \asymp \frac{\lambda_{\bd}(L_{h_j}^2)}{\lambda_{\bd}(L_{h_j})^2} h_j^{-d_j} L_{h_j}^2 (\bx^{(j)}, \by^{(j)}) f_j(\by^{(j)}).
\end{align*}

Note that
\begin{align*}
    \prod_{j=1}^{r} \mathsf{E} \lrb{ L_{h_j}^2 (\bx^{(j)}, \bX_{1}^{(j)})
    L_{h_j}^2 (\by^{(j)}, \bX_{1}^{(j)}) } & \asymp \lrb{ \frac{\lambda_{\bd}(L_{h_j}^2)}{\lambda_{\bd}(L_{h_j})^2} }^r \rho(\bh)^{-1} L_{\bh}^2 (\bx, \by) f_0(\by).
\end{align*}

\section{New proof - convergence}

\begin{align*}
    H_n(\bX_1, \bX_2) &= \int_{\Sr} \bar{L}_{\bh} (\bx,\bX_1) \bar{L}_{\bh} (\bx,\bX_2) \,\rd \bx.
\end{align*}

\begin{align*}
    \mathsf{E} ( H_n^4(\bX_1, \bX_2)) &= \mathsf{E}\lrb{ \lrb{\int_{\Sr} \bar{L}_{\bh} (\bx,\bX_1) \bar{L}_{\bh} (\bx,\bX_2) \,\rd \bx}^4 } \\
    &= \mathsf{E}\lrb{ \prod_{k=1}^{4} \lrb{\int_{\Sr} \bar{L}_{\bh} (\bx_k,\bX_1) \bar{L}_{\bh} (\bx_k,\bX_2) \,\rd \bx_k} } \\ 
    &= \mathsf{E} \lrb{\int_{(\Sr)^4} \prod_{k=1}^{4}\bar{L}_{\bh} (\bx_k,\bX_1) \bar{L}_{\bh} (\bx_k,\bX_2) \prod_{k=1}^{4} \,\rd \bx_k}  \\
    &= \int_{(\Sr)^4} \mathsf{E} \lrb{ \prod_{k=1}^{4}\bar{L}_{\bh} (\bx_k,\bX_1) \bar{L}_{\bh} (\bx_k,\bX_2) }\prod_{k=1}^{4} \,\rd \bx_k  \\
    &= \int_{(\Sr)^4} \mathsf{E} \lrb{ \prod_{k=1}^{4}\bar{L}_{\bh} (\bx_k,\bX_1)} \mathsf{E} \lrb{ \prod_{k=1}^{4} \bar{L}_{\bh} (\bx_k,\bX_2) }\prod_{k=1}^{4} \,\rd \bx_k  \\
    &= \int_{(\Sr)^4} \mathsf{E}^2 \lrb{ \prod_{k=1}^{4}\bar{L}_{\bh} (\bx_k,\bX_1)} \prod_{k=1}^{4} \,\rd \bx_k.
\end{align*}

Now note that
\begin{align*}
    \mathsf{E} \lrb{ \prod_{k=1}^{4}\bar{L}_{\bh} (\bx_k,\bX_1)} =&\; \mathsf{E} \lrb{ \prod_{k=1}^{4} \lrb{ L_{\bh} (\bx_k,\bX_1) - \mathsf{E} (L_{\bh} (\bx_k,\bX_1)) } } \\
    =&\; \mathsf{E} \lrb{ L_{\bh} (\bx_1,\bX_1) L_{\bh} (\bx_2,\bX_1) L_{\bh} (\bx_3,\bX_1) L_{\bh} (\bx_4,\bX_1)}  \\
    &- \mathsf{E} \lrb{ L_{\bh} (\bx_1,\bX_1) L_{\bh} (\bx_2,\bX_1) L_{\bh} (\bx_3,\bX_1) \mathsf{E} (L_{\bh} (\bx_4,\bX_1))} \\
    &+ \dots
\end{align*}
From $C_p$ inequality, it follows that the dominant terms of $\mathsf{E}^2 \lrb{ \prod_{k=1}^{4}\bar{L}_{\bh} (\bx_k,\bX_1)}$ will be the squares of the expectations of the terms in the decomposition above.

If we consider the first term, we have
\begin{align*}
    t_1 &\defin \int_{(\Sr)^4} \mathsf{E}^2 \lrb{ L_{\bh} (\bx_1,\bX_1) L_{\bh} (\bx_2,\bX_1) L_{\bh} (\bx_3,\bX_1) L_{\bh} (\bx_4,\bX_1)} \prod_{k=1}^{4} \,\rd \bx_k \\
    &= \int_{(\Sr)^4} \lrb{ \int_{\Sr}L_{\bh} (\bx_1,\by) L_{\bh} (\bx_2,\by) L_{\bh} (\bx_3,\by) L_{\bh} (\bx_4,\by) f(\by) \,\rd \by }^2 \prod_{k=1}^{4} \,\rd \bx_k.
\end{align*}

Now consider the tangent-normal decomposition $\bx_k^{(j)} = s_k^{(j)} \by_k^{(j)} + \lrb{1 - \lrb{s_k^{(j)}}^2}^{1/2} \bB_{\by^{(j)}} \bxi_k^{(j)}$, for $k = 1,2,3,4$ and $j =1,\dots,r$, where $\bxi_k^{(j)} \in \Sjm$, $\bB_{\by^{(j)}}$  is a $(d_j+1) \times d_j$ matrix such that $\bB_{\by^{(j)}}^{\top} \bB_{\by^{(j)}} = \bI_{d_j}$ and $\bB_{\by^{(j)}} \bB_{\by^{(j)}}^{\top} = \bI_{d_j+1} - \by^{(j)} \by^{(j) \top}$, and $s_k^{(j)} \in [-1,1]$. In such case, $\,\rd \bx_k^{(j)} = \lrb{1 - \lrb{s_k^{(j)}}^2}^{d_j/2 -1} \,\rd s_k^{(j)} \,\rd \bxi_k^{(j)}$. Given an arbitrary function $g$, this transformation of variables entails that
\begin{align*}
    \int_{\Sj} g(\bx^{(j)}) \,\rd \bx^{(j)} &= \int_{-1}^{1} \int_{\Sjm} g(s^{(j)} \by^{(j)} + (1-(s^{(j)})^2)^{1/2} \bB_{\by^{(j)}} \bxi^{(j)}) (1-(s^{(j)})^2)^{d_j/2 -1} \,\rd \bxi^{(j)} \,\rd s^{(j)}\\
    &= \int_{-1}^{1} \int_{\Sjm} g(s^{(j)} \by^{(j)} + \ab_{s^{(j)},\bxi^{(j)}}) J(s^{(j)}) \,\rd \bxi^{(j)} \,\rd s^{(j)},
\end{align*}
where $J(s^{(j)}) \defin \lrb{1 - (s^{(j)})^2}^{d_j/2 -1}$ and $\ab_{s^{(j)},\bxi^{(j)}}\defin (1-(s^{(j)})^2)^{1/2} \bB_{\by^{(j)}} \bxi^{(j)}$. Applying this tangent-normal $r$ times on $\Sr$, we have
\begin{align}
    \int_{\Sr} g(\bx) \,\rd \bx = \int_{[-1,1]^r} \int_{\Srm} g(s^{(1)} \by^{(1)} + \ab_{s^{(1)},\bxi^{(1)}},\ldots, s^{(r)} \by^{(r)} + \ab_{s^{(r)},\bxi^{(r)}}) \prod_{j=1}^r J(s^{(j)}) \,\rd \bs\,\rd \bxi.\label{eq:tnpoly}
\end{align}


\begin{align}
    \int_{\Sr} g(\bx) \,\rd \bx = \int_{\times_{j=1}^r[0,2h_j^{-2}]} \int_{\Srm} g((1-h^2_1s^{(1)}) \by^{(1)} + \ba_{\bx^{(1)},\bxi^{(1)},h},\ldots, (1-h^2_r s^{(r)}) \by^{(r)} + \ba_{\bx^{(1)},\bxi^{(1)},h}) \prod_{j=1}^r J_{h_j}(s^{(j)}) \,\rd \bs\,\rd \bxi.\label{eq:tnpoly}
\end{align}


To simplify the notation, consider $J_k(\bs_k) \defin \prod_{j=1}^r \lrb{1 - (\bs_k^{(j)})^2}^{d_j/2 -1}$ for $\bs_k\in[-1,1]^r$. Applying the tangent-normal decomposition in the variables $\bx_2,\bx_3,\bx_4$ in $t_1$, we obtain
\begin{align*}
    t_1 =&\; \int_{\Sr} \int_{([-1,1]^r)^3} \int_{(\Srm)^3} \lrb{ \int_{\Sr}L_{\bh} (\bx_1,\by) \prod_{k=2}^{4} c_{h,d}(L) L \lrb{ \frac{1 - s_k}{\bh^2} } f(\by) \,\rd \by }^2 \prod_{k=2}^{4} J_k(\bs_k) \,\rd \bxi \,\rd s \,\rd \bx_1.
\end{align*}
We can apply again the tangent-normal decomposition, this time on $\by^{(j)}$: $\by_k^{(j)} = s_k^{(j)} \bx_1^{(j)} + \lrb{1 - \lrb{s_k^{(j)}}^2}^{1/2} B_{x_1^{(j)}} \bxi_1^{(j)}$, for $j \in \{ 1, \dots, r \}$. This leads us to
\begin{align*}
    t_1 =&\; [c_{h,d}(L)]^8 \int_{\Sr} \int_{([-1,1]^r)^3} \int_{(\Srm)^3} \lrb{ \int_{[-1,1]^r} \int_{\Sr} \prod_{k=1}^{4} L \lrb{ \frac{1 - s_k}{\bh^2} } f \lrb{ s_1^{(j)} \bx_1^{(j)} + \lrb{1 - \lrb{s_1}^2}^{1/2} B_{x_1} \bxi_1 } \,\rd \bxi_1 J_1(s_1) \,\rd s_1 }^2 \\
    & \prod_{k=2}^{4} J_k(s_k) \,\rd \bxi \,\rd s \,\rd \bx_1.
\end{align*}

Now consider the transformation $t_k^{(j)} = \frac{1 - s_k^{(j)}}{h_j^2}$. Therefore, $\,\rd s_k^{(j)} = (-h_j^2) \,\rd t_k^{(j)}$ and
\begin{align*}
    t_1 =&\; [c_{h,d}(L)]^8 \int_{\Sr} \int_{I_h} \int_{(\Srm)^3} \lrb{ \int_{[-1,1]^r} \int_{\Sr} \prod_{k=1}^{4} L \lrb{ t_k} f \lrb{ x_1 + \alpha_{x_1,\bxi_1,h} } \rho (\bh) \prod_{j=1}^r (t_1^{(j)})^{d_j/2 -1} (2 -h_j^2 t_1^{(j)})^{d_j/2 -1} \,\rd \bxi_1 \,\rd t_1 }^2 \\
    & \prod_{k=2}^{4} J_k(t_k) \rho (\bh)^3 \,\rd \bxi \,\rd s \,\rd \bx_1 \\
    =&\; [c_{h,d}(L)]^8 \rho (\bh)^5 \omega_{\bd - 1}^3 \prod_{k=2}^4 \int_{I_h} L^4(t) J(t) \,\rd t \int_{\Sr} \lrb{ \int_{I_h} \int_{\Srm} L(t_1)  f \lrb{ x_1 + \alpha_{x_1,\bxi_1,h} } J_1(t_1) \,\rd \bxi_1 \,\rd t_1 }^2 \,\rd x_1.
\end{align*}
where $I_h \defin \times_{j=1}^r [0, 2h_j^{-2}]$.

Therefore, $t_1 = O([c_{h,d}(L)]^8 \rho (\bh)^5)$. Now remember that $c_{\bd,L}(\bh)^{-1} = \rho(\bh) \lambda_{\bh,\bd}(L)$, and, as a consequence, $c_{\bd,L}(\bh) \asymp \rho(\bh)^{-1}$. Therefore,
\begin{align*}
    t_1 &= O(\rho (\bh)^{-8} \rho (\bh)^5) \\
    &= O(\rho (\bh)^{-3}).
\end{align*}

Therefore, $\mathsf{E}[H_n^4(\bX_1,\bX_2)]=O(\rho(\bh)^{-3})$

%Now remember that $L_{h_j} (\bx^{(j)},\by^{(j)}) \defin c_{h,d_j}(L) L \Big( \frac{1 - \bx^{(j)\top} \by^{(j)}}{h_j^2} \Big)$. Replacing it, we have $O(\rho (\bh)^{-3}\rho (\bh)^{8}) = O(\rho (\bh)^{5})$.

\subsection{Second term}

Without loss of generality, consider the second term
\begin{align*}
    t_2 &\defin \int_{(\Sr)^4} \mathsf{E}^2 \lrb{ L_{\bh} (\bx_1,\bX_1) L_{\bh} (\bx_2,\bX_1) L_{\bh} (\bx_3,\bX_1) \mathsf{E} \lrb{ L_{\bh} (\bx_4,\bX_1)} }\prod_{k=1}^{4} \,\rd \bx_k \\
    &= \int_{(\Sr)^4} \mathsf{E}^2 \lrb{ L_{\bh} (\bx_4,\bX_1)} \mathsf{E}^2 \lrb{ L_{\bh} (\bx_1,\bX_1) L_{\bh} (\bx_2,\bX_1) L_{\bh} (\bx_3,\bX_1) }\prod_{k=1}^{4} \,\rd \bx_k \\
    &= \lrb{ \int_{\Sr} \mathsf{E}^2 \lrb{ L_{\bh} (\bx_4,\bX_1)} \,\rd \bx_4} \lrb{ \int_{(\Sr)^3} \mathsf{E}^2 \lrb{ L_{\bh} (\bx_1,\bX_1) L_{\bh} (\bx_2,\bX_1) L_{\bh} (\bx_3,\bX_1) }\prod_{k=1}^{3} \,\rd \bx_k } \\
    &= \lrb{ \int_{\Sr} \lrb{ \int_{\Sr} L_{\bh} (\bx_4,\by) f(\by) \,\rd \by }^2 \,\rd \bx_4} \lrb{ \int_{(\Sr)^3} \lrb{ \int_{\Sr} L_{\bh} (\bx_1,\by) L_{\bh} (\bx_2,\by) L_{\bh} (\bx_3,\by) f(\by) \,\rd \by}^2 \prod_{k=1}^{3} \,\rd \bx_k } \\
    & \defin \theta_{4} \times \theta_{123}.
    %&= \int_{(\Sr)^4} \lrb{ \int_{\Sr}L_{\bh} (\bx_1,\by) L_{\bh} (\bx_2,\by) L_{\bh} (\bx_3,\by) L_{\bh} (\bx_4,\by) f(\by) \,\rd \by }^2 \prod_{k=1}^{4} \,\rd \bx_k.
\end{align*}
Now, note that we can apply on both $\theta_{4}$ and $\theta_{123}$ calculations analogous to the calculations done from $t_1$. Let's consider $\theta_{123}$ first; applying tangent-normal transformation of variables, we obtain
\begin{align*}
    \theta_{123} &= \int_{\Sr} \int_{([-1,1]^r)^2} \int_{(\Srm)^2} \lrb{ \int_{\Sr}L_{\bh} (\bx_1,\by) c_{h,d}(L)^2 L \lrb{ \frac{1 - s_2}{\bh^2} } L \lrb{ \frac{1 - s_3}{\bh^2} } f(\by) \,\rd \by }^2 J_2(s_2) J_3(s_3) \,\rd \bxi \,\rd s \,\rd \bx_1 \\
    &= [c_{h,d}(L)]^6 \int_{\Sr} \int_{([-1,1]^r)^2} \int_{(\Srm)^2} \lrb{ \int_{[-1,1]^r} \int_{\Sr} \prod_{k=1}^{3} L \lrb{ \frac{1 - s_k}{\bh^2} } f \lrb{ s_1^{(j)} \bx_1^{(j)} + \lrb{1 - \lrb{s_1}^2}^{1/2} B_{x_1} \bxi_1 } \,\rd \bxi_1 J_1(s_1) \,\rd s_1 }^2 \\
    & J_2(s_2) J_3(s_3) \,\rd \bxi \,\rd s \,\rd \bx_1.
\end{align*}

\subsection{Third term}

Without loss of generality, consider the third term:
\begin{align*}
    t_3 &\defin \int_{(\Sr)^4} \mathsf{E}^2 \lrb{ L_{\bh} (\bx_1,\bX_1) L_{\bh} (\bx_2,\bX_1) \mathsf{E} \lrb{ L_{\bh} (\bx_3,\bX_1)} \mathsf{E} \lrb{ L_{\bh} (\bx_4,\bX_1)} }\prod_{k=1}^{4} \,\rd \bx_k \\
    &= \int_{(\Sr)^4} \mathsf{E}^2 \lrb{ L_{\bh} (\bx_1,\bX_1) L_{\bh} (\bx_2,\bX_1) } \mathsf{E}^2 \lrb{ L_{\bh} (\bx_3,\bX_1)} \mathsf{E}^2 \lrb{ L_{\bh} (\bx_4,\bX_1)} \prod_{k=1}^{4} \,\rd \bx_k \\
    &= \lrb{\int_{(\Sr)^2} \mathsf{E}^2 \lrb{ L_{\bh} (\bx_1,\bX_1) L_{\bh} (\bx_2,\bX_1) } \,\rd \bx_1 \,\rd \bx_2} \lrb{\int_{\Sr} \mathsf{E}^2 \lrb{ L_{\bh} (\bx_3,\bX_1)} \,\rd \bx_3} \lrb{\int_{\Sr} \mathsf{E}^2 \lrb{ L_{\bh} (\bx_4,\bX_1)} \,\rd \bx_4} \\
    &\defin \theta_{12} \times \theta_3 \times \theta_4
\end{align*}

\subsection{Fourth term}

Consider, finally, the last term:
\begin{align*}
    t_4 &\defin \int_{(\Sr)^4} \mathsf{E}^2 \lrb{ \mathsf{E} \lrb{ L_{\bh} (\bx_1,\bX_1)} \mathsf{E} \lrb{ L_{\bh} (\bx_2,\bX_1)} \mathsf{E} \lrb{ L_{\bh} (\bx_3,\bX_1)} \mathsf{E} \lrb{ L_{\bh} (\bx_4,\bX_1)} }\prod_{k=1}^{4} \,\rd \bx_k \\
    &= \int_{(\Sr)^4} \mathsf{E}^2 \lrb{ L_{\bh} (\bx_1,\bX_1) } \mathsf{E}^2 \lrb{L_{\bh} (\bx_2,\bX_1) } \mathsf{E}^2 \lrb{ L_{\bh} (\bx_3,\bX_1)} \mathsf{E}^2 \lrb{ L_{\bh} (\bx_4,\bX_1)} \prod_{k=1}^{4} \,\rd \bx_k \\
    &= \prod_{k=1}^{4} \lrb{\int_{\Sr}  \mathsf{E}^2 \lrb{ L_{\bh} (\bx_k,\bX_1)} \,\rd \bx_k} \\
    & \defin \theta_1 \theta_2 \theta_3 \theta_4.
\end{align*}

\section{Old calculations - convergence}

\textit{Asymptotic behavior of $\mathsf{E} (H_n^4(\bX_1, \bX_k))$.} Consider the following calculations.

\begin{align*}
    H_n(\bX_1, \bX_2) &= \frac{2}{n^2} \int_{\Sr} \bar{L}_{\bh} (\bx,\bX_1) \bar{L}_{\bh} (\bx,\bX_2) \,\rd \bx.
\end{align*}

\begin{align*}
    H_n^4(\bX_1, \bX_2) &= \frac{16}{n^{8}} \lrp{\int_{\Sr} \bar{L}_{\bh} (\bx,\bX_1) \bar{L}_{\bh} (\bx,\bX_2) \,\rd \bx}^4\\
    &=\frac{16}{n^{8}}\int_{(\Sr)^4} \Bigg( \prod_{l=1}^{4} \bar{L}_{\bh} (\bx_l,\bX_1) \Bigg) \Bigg( \prod_{l=1}^{4} \bar{L}_{\bh} (\bx_l,\bX_2) \Bigg) \prod_{l=1}^{4} \,\rd 
    \bx_l.
\end{align*}

\begin{align*}
    \mathsf{E} (H_n^4(\bX_1, \bX_2)) &= \frac{16}{n^{8}} \int_{(\Sr)^4} \mathsf{E}^2 \Bigg( \prod_{l=1}^{4} \bar{L}_{\bh} (\bx_l,\bX_1) \Bigg) \prod_{l=1}^{4} \,\rd 
    \bx_l \\
    & \le \frac{16}{n^{8}} \int_{(\Sr)^4} \mathsf{E} \Bigg( \prod_{l=1}^{4} \bar{L}_{\bh}^2 (\bx_l,\bX_1) \Bigg) \prod_{l=1}^{4} \,\rd \bx_l \\
    &= \frac{16}{n^{8}} \mathsf{E} \Bigg[ \int_{(\Sr)^4} \Bigg( \prod_{l=1}^{4} \bar{L}_{\bh}^2 (\bx_l,\bX_1) \Bigg) \prod_{l=1}^{4} \,\rd \bx_l \Bigg] \\
    &= \frac{16}{n^{8}} \mathsf{E} \Bigg[ \Bigg( \int_{\Sr} \bar{L}_{\bh}^2 (\bx,\bX_1) \,\rd \bx \Bigg)^4 \Bigg] \\
    &= \frac{16}{n^{8}} \mathsf{E} (Z_{n1}^4).
\end{align*}

Applying the $c_p$ inequality two times, we have that
\begin{align*}
    Z_{n1}^4 & \le \Bigg( \int_{\Sr} 2[L_{\bh}^2 (\bx,\bX_1) + \mathsf{E}^2 (L_{\bh} (\bx,\bX_1))] \rd \bx \Bigg)^4 \\
    &= 16 \Bigg( \int_{\Sr} L_{\bh}^2 (\bx,\bX_1) \rd \bx + \int_{\Sr} \mathsf{E}^2 (L_{\bh} (\bx,\bX_1)) \rd \bx \Bigg)^4 \\
    & \le 128 \Bigg( \Bigg( \int_{\Sr} L_{\bh}^2 (\bx,\bX_1) \rd \bx \Bigg)^4 + \Bigg( \int_{\Sr} \mathsf{E}^2 (L_{\bh} (\bx,\bX_1)) \rd \bx \Bigg)^4 \Bigg).
\end{align*}

Now applying Lemma \ref{lemma-equivLh}, 
\begin{align*}
    \int_{\Sr} L_{\bh}^2 (\bx,\bX_1) \rd \bx \asymp \frac{\lambda_{\bd}(L^2)}{\lambda_{\bd}(L)^2} \rho(\bh)^{-1}.
\end{align*}
\begin{align*}
    \int_{\Sr} \mathsf{E}^2 (L_{\bh} (\bx,\bX_i)) &= \int_{\Sr} \Bigg( \int_{\Sr} L_{\bh} (\bx,\bX_1) f(\bx) \rd \bx \Bigg)^2 \rd \bx \\
    & \asymp \int_{\Sr} \Bigg( \frac{\lambda_{\bd}(L)}{\lambda_{\bd}(L)} f(\bx) \Bigg)^2 \rd \bx \\
    &= R(f).
\end{align*}

Therefore,
\begin{align*}
    \mathsf{E} (H_n^4(\bX_1, \bX_k)) &= \frac{16}{n^{8}} O(\rho(\bh)^{-4} + R(f)^4) \\
    &= O(n^{-8} \rho(\bh)^{-4}).
\end{align*}

\textit{Asymptotic behavior of $\mathsf{E} (G_n^2(\bX_1,\bX_2))$.} Consider the following calculations.

\begin{align*}
    H_n(\bx,\bX_1) &= \frac{2}{n^2} \int_{\Sr} \bar{L}_{\bh} (\bt,\bx) \bar{L}_{\bh} (\bt,\bX_1) \,\rd \bt.
\end{align*}

\begin{align*}
    G_n(\bx,\by) &= \mathsf{E} (H_n(\bX_1,\bx)H_n(\bX_1,\by))\\
    &=\int_{\Sr} \lrb{\frac{2}{n^2} \int_{\Sr} \bar{L}_{\bh} (\bt,\bx) \bar{L}_{\bh} (\bt,\bx_1) \,\rd \bt} \lrb{\frac{2}{n^2} \int_{\Sr} \bar{L}_{\bh} (\bs,\by) \bar{L}_{\bh} (\bs,\bx_1) \,\rd \bs} f(\bx_1) \,\rd \bx_1 \\
    &= \frac{4}{n^4} \int_{\Sr} \int_{\Sr} \int_{\Sr} \bar{L}_{\bh} (\bt,\bx) \bar{L}_{\bh} (\bt,\bx_1) \bar{L}_{\bh} (\bs,\by) \bar{L}_{\bh} (\bs,\bx_1) f(\bx_1) \,\rd \bx_1 \,\rd \bt \,\rd \bs \\
    &= \frac{4}{n^4} \int_{\Sr} \int_{\Sr} \bar{L}_{\bh} (\bt,\bx) \bar{L}_{\bh} (\bs,\by) \mathsf{E} (\bar{L}_{\bh} (\bt,\bX_1) \bar{L}_{\bh} (\bs,\bX_1)) \,\rd \bt \,\rd \bs.
\end{align*}

Note that, by Cauchy--Schwarz inequality,
\begin{align*}
    \mathsf{E} (\bar{L}_{\bh} (\bt,\bX_1) \bar{L}_{\bh} (\bs,\bX_1)) & \le \sqrt{\mathsf{E}^2 (\bar{L}_{\bh} (\bt,\bX_1)) \mathsf{E}^2 (\bar{L}_{\bh} (\bs,\bX_1))} \\
    &= \sqrt{\mathsf{Var} ({L}_{\bh} (\bt,\bX_1)) \mathsf{Var} ({L}_{\bh} (\bs,\bX_1))}.
\end{align*}

Therefore,
\begin{align*}
    G_n(\bx,\by) & \le \frac{4}{n^4} \int_{(\Sr)^2}  \bar{L}_{\bh} (\bt,\bx) \bar{L}_{\bh} (\bs,\by) \sqrt{\mathsf{Var} ({L}_{\bh} (\bt,\bX_1)) \mathsf{Var} ({L}_{\bh} (\bs,\bX_1))} \,\rd \bt \,\rd \bs. \\
    G_n^2(\bx,\by) & \le \frac{16}{n^{8}} \int_{(\Sr)^4} \bar{L}_{\bh} (\bt_1,\bx) \bar{L}_{\bh} (\bt_2,\by) \bar{L}_{\bh} (\bt_3,\bx) \bar{L}_{\bh} (\bt_4,\by) \sqrt{\prod_{i=1}^{4} \mathsf{Var} ({L}_{\bh} (\bt_i,\bX_1)) } \prod_{i=1}^{4} \,\rd \bt_i.
\end{align*}

Therefore,
\begin{align*}
    \mathsf{E} (G_n^2(\bX_1,\bX_2)) & \le \frac{16}{n^{8}} \int_{(\Sr)^4} \mathsf{E} \lrb{ \bar{L}_{\bh} (\bt_1,\bX_1) \bar{L}_{\bh} (\bt_2,\bX_2) \bar{L}_{\bh} (\bt_3,\bX_1) \bar{L}_{\bh} (\bt_4,\bX_2)} \sqrt{\prod_{i=1}^{4} \mathsf{Var} ({L}_{\bh} (\bt_i,\bX_1)) } \prod_{i=1}^{4} \,\rd \bt_i \\
    &= \frac{16}{n^{8}} \int_{(\Sr)^4} \mathsf{E} \lrb{ \bar{L}_{\bh} (\bt_1,\bX_1) \bar{L}_{\bh} (\bt_3,\bX_1) } \mathsf{E} \lrb{ \bar{L}_{\bh} (\bt_2,\bX_2) \bar{L}_{\bh} (\bt_4,\bX_2)} \sqrt{\prod_{i=1}^{4} \mathsf{Var} ({L}_{\bh} (\bt_i,\bX_1)) } \prod_{i=1}^{4} \,\rd \bt_i \\
    &= \frac{16}{n^{8}} \int_{(\Sr)^4} \mathsf{E} \lrb{ \bar{L}_{\bh} (\bt_1,\bX_1) \bar{L}_{\bh} (\bt_3,\bX_1) } \mathsf{E} \lrb{ \bar{L}_{\bh} (\bt_2,\bX_1) \bar{L}_{\bh} (\bt_4,\bX_1)} \sqrt{\prod_{i=1}^{4} \mathsf{Var} ({L}_{\bh} (\bt_i,\bX_1)) } \prod_{i=1}^{4} \,\rd \bt_i. \\
\end{align*}

Applying Cauchy--Schwartz inequality again, we obtain
\begin{align*}
    \mathsf{E} (G_n^2(\bX_1,\bX_2)) \le &\; \frac{16}{n^{8}} \int_{(\Sr)^4}  \sqrt{\mathsf{Var} ({L}_{\bh} (\bt_1,\bX_1)) \mathsf{Var} ({L}_{\bh} (\bt_3,\bX_1))} \sqrt{\mathsf{Var} ({L}_{\bh} (\bt_2,\bX_1)) \mathsf{Var} ({L}_{\bh} (\bt_4,\bX_1))} \\
    & \times \sqrt{\prod_{i=1}^{4} \mathsf{Var} ({L}_{\bh} (\bt_i,\bX_1)) } \prod_{i=1}^{4} \,\rd \bt_i \\
    =&\; \frac{16}{n^{8}} \int_{(\Sr)^4} \sqrt{\prod_{i=1}^{4} \mathsf{Var} ({L}_{\bh} (\bt_i,\bX_1)) } \sqrt{\prod_{i=1}^{4} \mathsf{Var} ({L}_{\bh} (\bt_i,\bX_1)) } \prod_{i=1}^{4} \,\rd \bt_i \\
    =&\; \frac{16}{n^{8}} \int_{(\Sr)^4} \prod_{i=1}^{4} \mathsf{Var} ({L}_{\bh} (\bt_i,\bX_1)) \prod_{i=1}^{4} \,\rd \bt_i \\
    =&\; \frac{16}{n^{8}} \int_{(\Sr)^4} \prod_{i=1}^{4} n \mathsf{Var} (\hat{f}(\bt_i;\bh)) \prod_{i=1}^{4} \,\rd \bt_i \\
    =&\; \frac{16}{n^{8}} \int_{(\Sr)^4} \prod_{i=1}^{4} n \lrb{\frac{v_{\bd}(L)}{n \rho(\bh)} f(\bt_i)(1 + o(1))} \prod_{i=1}^{4} \,\rd \bt_i \\
    =&\; \frac{16}{n^{8}} \int_{(\Sr)^4} \prod_{i=1}^{4} \frac{v_{\bd}(L)}{ \rho(\bh)} f(\bt_i)(1 + o(1)) \prod_{i=1}^{4} \,\rd \bt_i \\
    =&\; \frac{16}{n^{8}} \prod_{i=1}^{4} \int_{(\Sr)^4} \frac{v_{\bd}(L)}{ \rho(\bh)} f(\bt_i)(1 + o(1)) \,\rd \bt_i \\
    =&\; \frac{16}{n^{8}} \prod_{i=1}^{4} \frac{v_{\bd}(L)}{ \rho(\bh)} (1 + o(1)) \\
    =&\; \frac{16}{n^{8}} \frac{v_{\bd}^4(L)}{ \rho^4(\bh)} (1 + o(1)).
\end{align*}

As a consequence,
\begin{align*}
    \mathsf{E} (G_n^2(\bX_1,\bX_2)) &= O(n^{-8} \rho^{-4}(\bh)).
\end{align*}

\textit{Asymptotic behavior of the ratio $A_n / B_n^2$.} Defining $A_n \defin n^3 \mathsf{E} (H_n^4(\bX_1, \bX_2)) + n^4 \mathsf{E} (G_n^2(\bX_1,\bX_2))$, then, from the previous results,
\begin{align*}
    A_n &= O(n^3 n^{-8}\rho(\bh)^{-4} + n^4 n^{-8}\rho(\bh)^{-4}) \\
    &= O(n^{-5} \rho^{-4}(\bh)+n^{-4} \rho^{-4}(\bh))\\
    &= O(n^{-4} \rho^{-4}(\bh)).
\end{align*}

\begin{align*}
    B_n &= \mathsf{E}[H_n^2(\bX_1,\bX_2)]=\frac{2}{n^2} \int_{\Sr} \int_{\Sr} \mathsf{E}^2 ( \bar{L}_{\bh} (\bx,\bX_1) \bar{L}_{\bh} (\by,\bX_1) ) \,\rd \bx \,\rd \by
\end{align*}

\begin{comment}
Note that
\begin{align*}
    \mathsf{E} ( \bar{L}_{\bh} (\bx,\bX_1) \bar{L}_{\bh} (\by,\bX_1) ) &= \mathsf{E} ( L_{\bh} (\bx,\bX_1) L_{\bh} (\by,\bX_1) ) - \mathsf{E} (L_{\bh} (\bx,\bX_1)) \mathsf{E} (L_{\bh} (\by,\bX_1)).
\end{align*}

Furthermore,
\begin{align*}
    \mathsf{E} ( L_{\bh} (\bx,\bX_1) L_{\bh} (\by,\bX_1) ) &= \int_{\Sr} L_{\bh} (\bx,\bt) L_{\bh} (\by,\bt) f(\bt) \,\rd \bt \\
    & \asymp L_{\bh} (\bx,\by) f(\by).\\
    \mathsf{E} (L_{\bh} (\bx,\bX_1)) &= \int_{\Sr} L_{\bh} (\bx,\bt) f(\bt) \,\rd \bt \\
    & \asymp f(\bx).\\
    \mathsf{E} (L_{\bh} (\by,\bX_1)) &= \int_{\Sr} L_{\bh} (\by,\bt) f(\bt) \,\rd \bt \\
    & \asymp f(\by).\\
\end{align*}

As a consequence,
\begin{align*}
    \mathsf{E}^2 ( \bar{L}_{\bh} (\bx,\bX_1) \bar{L}_{\bh} (\by,\bX_1) ) & \asymp \lrb{L_{\bh} (\bx,\by) f(\by) - f(\bx)f(\by)}^2 \\
    &= L_{\bh}^2 (\bx,\by) f^2(\by) + f^2(\bx)f^2(\by) -2L_{\bh} (\bx,\by) f(\bx) f^2(\by).
\end{align*}

Let's integrate each term.
\begin{align*}
    \int_{\Sr} \int_{\Sr} L_{\bh}^2 (\bx,\by) f^2(\by) \,\rd \by \,\rd \bx & \asymp \int_{\Sr} \frac{\lambda_{\bd}(L^2)}{\lambda_{\bd}(L)^2} \rho(\bh)^{-1} f^2(\bx) \,\rd \bx \\
    &= \frac{\lambda_{\bd}(L^2)}{\lambda_{\bd}(L)^2} \rho(\bh)^{-1} R(f). \\
    \int_{\Sr} \int_{\Sr} f^2(\bx)f^2(\by) \,\rd \by \,\rd \bx &= R^2(f). \\
    \int_{\Sr} \int_{\Sr}-2L_{\bh} (\bx,\by) f(\bx) f^2(\by) \,\rd \by \,\rd \bx & \asymp -2 \int_{\Sr} f(\bx) f^2(\bx) \,\rd \bx\\
    &= -2 \int_{\Sr}f^3(\bx) \,\rd \bx.
\end{align*}

Finally,
\begin{align*}
    B_n \asymp \frac{2}{n^2} \lrb{\frac{\lambda_{\bd}(L^2)}{\lambda_{\bd}(L)^2} \rho(\bh)^{-1} R(f) + R^2(f) -2 \int_{\Sr}f^3(\bx) \,\rd \bx}.
\end{align*}
\end{comment}

$B_n\asymp 4 n^{-4}\rho(\bh)^{-1} \sigma^2$


\begin{align*}
    \frac{\mathsf{E}[G_n^2(\bX_1,\bX_2)]+n^{-1}\mathsf{E}[H_n^4(\bX_1,\bX_2)]}{\lrb{\mathsf{E}[H_n^2(\bX_1,\bX_2)]}^2}&\asymp \frac{\mathsf{E}[G_n^2(\bX_1,\bX_2)]+n^{-1}\mathsf{E}[H_n^4(\bX_1,\bX_2)]}{\lrb{4 n^{-4}\rho(\bh)^{-1} \sigma^2}^2}\\
    &=O\lrp{\frac{\mathsf{E}[G_n^2(\bX_1,\bX_2)]+n^{-1}\mathsf{E}[H_n^4(\bX_1,\bX_2)]}{n^{-8}\rho(\bh)^{-2}}}\\
    &=O\lrp{\frac{n^{-8}\rho(\bh)^{-4}+n^{-1} n^{-8} \rho(\bh)^{-4}}{n^{-8}\rho(\bh)^{-2}}}\\
    &=O\lrp{\rho(\bh)^{-2}+n^{-1}  \rho(\bh)^{-2}}\\
\end{align*}

Notice that, therefore, $A_n = O(n^{-12} \rho^{-4}(\bh))$ and $B_n^2 \asymp n^{-4}(\rho(\bh)^{-2}+1)$, which entails that $A_n / B_n^2 \to 0$ as $n$ goes to infinity.

\section{Lemmas}

\begin{lemma}\label{lemma-tnd-poly}
    Let $\bx \in \Sr$ a random vector supported in the polysphere. Let $\bxi_k^{(j)} \in \Sjm$, $\bB_{\by^{(j)}}$  is a $(d_j+1) \times d_j$ matrix such that $\bB_{\by^{(j)}}^{\top} \bB_{\by^{(j)}} = \bI_{d_j}$ and $\bB_{\by^{(j)}} \bB_{\by^{(j)}}^{\top} = \bI_{d_j+1} - \by^{(j)} \by^{(j) \top}$, and $s_k^{(j)} \in [-1,1]$. Given an arbitrary function $g$, it is true that
    \begin{align}
    \int_{\Sr} g(\bx) \,\rd \bx = \int_{\Srm} \int_{[-1,1]^r} g(s^{(1)} \by^{(1)} + \ab_{s^{(1)},\bxi^{(1)}},\ldots, s^{(r)} \by^{(r)} + \ab_{s^{(r)},\bxi^{(r)}}) \prod_{j=1}^r J(s^{(j)}) \,\rd \bs\,\rd \bxi,\label{eq:tnpoly}
\end{align}
where $J(s^{(j)}) \defin \lrb{1 - (s^{(j)})^2}^{d_j/2 -1}$ and $\ab_{s^{(j)},\bxi^{(j)}}\defin (1-(s^{(j)})^2)^{1/2} \bB_{\by^{(j)}} \bxi^{(j)}$. 
\end{lemma}

\begin{lemma}\label{lemma-tnd2-poly}
    Under the same definitions of the Lemma \ref{lemma-tnd-poly},
    \begin{align}
    \int_{\Sr} g(\bx) \,\rd \bx = \int_{\times_{j=1}^r[0,2h_j^{-2}]} \int_{\Srm} g((1-h^2_1s^{(1)}) \by^{(1)} + \ba_{\bx^{(1)},\bxi^{(1)},h},\ldots, (1-h^2_r s^{(r)}) \by^{(r)} + \ba_{\bx^{(1)},\bxi^{(1)},h}) \prod_{j=1}^r J_{h_j}(s^{(j)}) \,\rd \bs\,\rd \bxi.\label{eq:tnpoly}
\end{align}
\end{lemma}

\section{Variance of $T_{n,2}$}

\begin{align*}
    T_{n,2} & \defin \int_{\Sr} [ \hat{f}_{0}(\bx;\bh) - \mathsf{E} (\hat{f}_{0}(\bx;\bh)) ]^2 \,\rd \bx \\
    &= \int_{\Sr} \lrb{ \prod_{j=1}^{r} \hat{f}_{j}(\bx^{(j)};h_j) - \mathsf{E} \lrb{\prod_{j=1}^{r} \hat{f}_{j}(\bx^{(j)};h_j)} }^2 \,\rd \bx \\
    &= \int_{\Sr} \lrb{ \prod_{j=1}^{r} \hat{f}_{j}(\bx^{(j)};h_j) - \prod_{j=1}^{r} \mathsf{E} \lrb{ \hat{f}_{j}(\bx^{(j)};h_j)} }^2 \,\rd \bx \\
    &= \int_{\Sr} \lrb{ \lrb{\prod_{j=1}^{r} \hat{f}_{j}(\bx^{(j)};h_j)}^2 + \lrb{\prod_{j=1}^{r} \mathsf{E} \lrb{ \hat{f}_{j}(\bx^{(j)};h_j)}}^2 - 2\lrb{\prod_{j=1}^{r} \hat{f}_{j}(\bx^{(j)};h_j) \prod_{j=1}^{r} \mathsf{E} \lrb{ \hat{f}_{j}(\bx^{(j)};h_j)}} } \,\rd \bx \\
    &= \int_{\Sr} \lrb{ \prod_{j=1}^{r} \hat{f}_{j}^2(\bx^{(j)};h_j) + \prod_{j=1}^{r} \mathsf{E}^2 \lrb{ \hat{f}_{j}(\bx^{(j)};h_j)} - 2\lrb{\prod_{j=1}^{r} \hat{f}_{j}(\bx^{(j)};h_j) \mathsf{E} \lrb{ \hat{f}_{j}(\bx^{(j)};h_j)}} } \,\rd \bx \\
    &= \int_{\Sr} \prod_{j=1}^{r} \hat{f}_{j}^2(\bx^{(j)};h_j) \,\rd \bx + \int_{\Sr} \prod_{j=1}^{r} \mathsf{E}^2 \lrb{ \hat{f}_{j}(\bx^{(j)};h_j)} \,\rd \bx - 2 \int_{\Sr} \prod_{j=1}^{r} \hat{f}_{j}(\bx^{(j)};h_j) \mathsf{E} \lrb{ \hat{f}_{j}(\bx^{(j)};h_j)}  \,\rd \bx \\
    &= \prod_{j=1}^{r} \int_{\Sr} \hat{f}_{j}^2(\bx^{(j)};h_j) \,\rd \bx^{(j)} + \prod_{j=1}^{r} \int_{\Sr} \mathsf{E}^2 \lrb{ \hat{f}_{j}(\bx^{(j)};h_j)} \,\rd \bx^{(j)} - 2 \prod_{j=1}^{r} \int_{\Sr} \hat{f}_{j}(\bx^{(j)};h_j) \mathsf{E} \lrb{ \hat{f}_{j}(\bx^{(j)};h_j)}  \,\rd \bx^{(j)}. \\
\end{align*}
Notice that $\prod_{j=1}^{r} \int_{\Sr} \mathsf{E}^2 \lrb{ \hat{f}_{j}(\bx^{(j)};h_j)} \,\rd \bx^{(j)}$ is a constant, therefore,
\begin{align*}
    \mathsf{Var}(T_{n,2}) &= \mathsf{Var} \lrb{\prod_{j=1}^{r} \int_{\Sr} \hat{f}_{j}^2(\bx^{(j)};h_j) \,\rd \bx^{(j)} - 2 \prod_{j=1}^{r} \int_{\Sr} \hat{f}_{j}(\bx^{(j)};h_j) \mathsf{E} \lrb{ \hat{f}_{j}(\bx^{(j)};h_j)}  \,\rd \bx^{(j)}} \\
    &= O \lrb{ \mathsf{Var} \lrb{\prod_{j=1}^{r} \int_{\Sr} \hat{f}_{j}^2(\bx^{(j)};h_j) \,\rd \bx^{(j)}} + \mathsf{Var} \lrb{- 2 \prod_{j=1}^{r} \int_{\Sr} \hat{f}_{j}(\bx^{(j)};h_j) \mathsf{E} \lrb{ \hat{f}_{j}(\bx^{(j)};h_j)}  \,\rd \bx^{(j)}}},
\end{align*}
where last equality comes from Lemma XXX. To simplify the notation, we can define $t_1 \defin \prod_{j=1}^{r} \int_{\Sr} \hat{f}_{j}^2(\bx^{(j)};h_j) \,\rd \bx^{(j)}$ and $t_2 \defin \prod_{j=1}^{r} \int_{\Sr} \hat{f}_{j}(\bx^{(j)};h_j) \mathsf{E} \lrb{ \hat{f}_{j}(\bx^{(j)};h_j)}  \,\rd \bx^{(j)}$, and, therefore,
\begin{align*}
    \mathsf{Var}(T_{n,2}) &= O(\mathsf{Var} (t_1)) + O(\mathsf{Var} (t_2)).
\end{align*}

Let's study first $O(\mathsf{Var} (t_1))$. We will compute the first two moments of $t_1$.
\begin{align*}
    \mathsf{E} (t_1) &= \mathsf{E} \lrb{\prod_{j=1}^{r} \int_{\Sr} \hat{f}_{j}^2(\bx^{(j)};h_j) \,\rd \bx^{(j)}} \\
    &= \prod_{j=1}^{r} \mathsf{E} \lrb{ \int_{\Sr} \hat{f}_{j}^2(\bx^{(j)};h_j) \,\rd \bx^{(j)}} \\
    &= \prod_{j=1}^{r} \int_{\Sr} \mathsf{E} \lrb{ \hat{f}_{j}^2(\bx^{(j)};h_j)} \,\rd \bx^{(j)}. \\
\end{align*}
Therefore,
\begin{align*}
    (\mathsf{E} (t_1))^2 &= \lrb{\prod_{j=1}^{r} \int_{\Sr} \mathsf{E} \lrb{ \hat{f}_{j}^2(\bx^{(j)};h_j)} \,\rd \bx^{(j)}} \lrb{\prod_{k=1}^{r} \int_{\Sr} \mathsf{E} \lrb{ \hat{f}_{k}^2(\by^{(j)};h_k)} \,\rd \by^{(k)}} \\
    &= \prod_{j=1}^{r} \lrb{\int_{\Sr} \mathsf{E} \lrb{ \hat{f}_{j}^2(\bx^{(j)};h_j)} \,\rd \bx^{(j)}} \lrb{\int_{\Sr} \mathsf{E} \lrb{ \hat{f}_{j}^2(\by^{(j)};h_j)} \,\rd \by^{(j)}} \\
     &= \prod_{j=1}^{r} \lrb{\int_{\Sr} \mathsf{E} \lrb{ \hat{f}_{j}^2(\bx^{(j)};h_j)} \,\rd \bx^{(j)}}^2.
\end{align*}

\begin{align*}
    \mathsf{E} (t_1^2) &= \mathsf{E} \lrb{ \lrb{\prod_{j=1}^{r} \int_{\Sr} \hat{f}_{j}^2(\bx^{(j)};h_j) \,\rd \bx^{(j)}}^2 }\\
    &= \mathsf{E} \lrb{ \prod_{j=1}^{r} \lrb{ \int_{\Sr} \hat{f}_{j}^2(\bx^{(j)};h_j) \,\rd \bx^{(j)}}^2 }\\
    &= \prod_{j=1}^{r} \mathsf{E} \lrb{ \lrb{ \int_{\Sr} \hat{f}_{j}^2(\bx^{(j)};h_j) \,\rd \bx^{(j)}}^2 }\\
    &= \prod_{j=1}^{r} \mathsf{E} \lrb{ \lrb{ \int_{\Sr} \hat{f}_{j}^2(\bx^{(j)};h_j) \,\rd \bx^{(j)}} \lrb{ \int_{\Sr} \hat{f}_{j}^2(\by^{(j)};h_j) \,\rd \by^{(j)}} }\\
    &= \prod_{j=1}^{r} \mathsf{E} \lrb{  \int_{\Sr} \int_{\Sr} \hat{f}_{j}^2(\bx^{(j)};h_j) \hat{f}_{j}^2(\by^{(j)};h_j) \,\rd \bx^{(j)} \,\rd \by^{(j)} }\\
    &= \prod_{j=1}^{r} \int_{\Sr} \int_{\Sr} \mathsf{E} \lrb{ \hat{f}_{j}^2(\bx^{(j)};h_j) \hat{f}_{j}^2(\by^{(j)};h_j) } \,\rd \bx^{(j)} \,\rd \by^{(j)}. \\
\end{align*}

%-------------------------------%
\section{Simulation study}
\label{sec:simulation}
%-------------------------------%

%Code to replicate plots in \cite{Garcia-Portugues2024}: \url{https://github.com/egarpor/polykde/tree/main/replication}

To assess the power of the test in different kinds of distributions, the following scenarios were considered.

\begin{enumerate}
    \item \textbf{Mixture of iid von Mises--Fisher distributions:} we simulate a latent random vector $\bmu \sim \mathrm{vMF}_d (\bmu_1, \kappa_1)$, and, then, we simulate $\bX^{(1)}, \dots, \bX^{(r)} | \bmu \stackrel {iid}\sim  \mathrm{vMF}_d (\bmu, \kappa)$. This produces the density
    \begin{align*}
        f(\bx) &= \int_{\Sr} c^\mathrm{vMF}_{d}(\kappa_1) e^{\kappa_1\bmu^{\top}\bmu_1} \prod_{j=1}^{r} c^\mathrm{vMF}_{d}(\kappa) e^{\kappa\bx^{(j) \top}\bmu} \,\rd \bmu \\
        &= c^\mathrm{vMF}_{d}(\kappa_1) c^\mathrm{vMF}_{d}(\kappa)^r \int_{\Sr} e^{\kappa_1\bmu^{\top}\bmu_1 + \sum_{j=1}^{r} \kappa\bx^{(j) \top}\bmu} \,\rd \bmu\\
        &= c^\mathrm{vMF}_{d}(\kappa_1) c^\mathrm{vMF}_{d}(\kappa)^r \int_{\Sr} e^{\|\kappa_1\bmu_1 + \sum_{j=1}^{r} \kappa\bx^{(j)}\|\bmu^{\top}(\kappa_1\bmu_1 + \sum_{j=1}^{r} \kappa\bx^{(j)})/\|\kappa_1\bmu_1 + \sum_{j=1}^{r} \kappa\bx^{(j)}\|} \,\rd \bmu\\
        &= \frac{c^\mathrm{vMF}_{d}(\kappa_1) c^\mathrm{vMF}_{d}(\kappa)^r}{ c^\mathrm{vMF}_{d}(\|\kappa_1\bmu_1 + \sum_{j=1}^{r} \kappa\bx^{(j)}\|)^r}.
\end{align*}
%{\color{red} Relation with \cite{Johnson1978,Wehrly1980}'s copula model for $d_1=d_2=1$ and $r=2$?}
It is worth noting that, for the case with $r=2, d_1 = d_2 = 1, \kappa_1 = 0$ (two random vectors supported in the circle), the joint distribution of the angles belongs to the  Johnson-Wehrly family of distribution - see \cite{Johnson1978,Wehrly1980}. This result is available in Proposition \ref{prop-js-family}.
    For convenience, we will reparametrize the model as $\delta \defin 1/\kappa_1$, because, in this way, for fixed parameters $\bmu_1$ and $\kappa$, the parameter $\delta$ will quantify the amount of deviation from independence. By convention, we will define that, when $\delta = 0$, then $\bmu = \bmu_1$ almost surely (degenerate random vector); therefore,
    \begin{align*}
        f(\bx) &= c^\mathrm{vMF}_{d}(1/\delta) c^\mathrm{vMF}_{d}(\kappa)^r \int_{\Sr} e^{(1/\delta)\bmu^{\top}\bmu_1 + \sum_{j=1}^{r} \kappa\bx^{(j) \top}\bmu} \,\rd \bmu,
    \end{align*}
    and, for $\delta = 0$,
    \begin{align*}
        f(\bx) &= \prod_{j=1}^{r} c^\mathrm{vMF}_{d}(\kappa) e^{\kappa\bx^{(j) \top}\bmu_1} \\
        &= c^\mathrm{vMF}_{d}(\kappa)^r e^{\sum_{j=1}^{r} \kappa\bx^{(j) \top}\bmu_1},
    \end{align*}
    which is the independence case.

    The following scenarios in the sphere are considered.
    \begin{enumerate}
        \item $r=3, d=2, \bmu_1 = (\sqrt{1/3})_{1 \times 9}, \kappa = 10, \delta = 0$ (independence case);
        \item $r=3, d=2, \bmu_1 = (\sqrt{1/3})_{1 \times 9}, \kappa = 10, \delta = 0.1$;
        \item $r=3, d=2, \bmu_1 = (\sqrt{1/3})_{1 \times 9}, \kappa = 10, \delta = 0.5$;
        \item $r=3, d=2, \bmu_1 = (\sqrt{1/3})_{1 \times 9}, \kappa = 10, \delta = 1$;
        \item $r=3, d=2, \bmu_1 = (\sqrt{1/3})_{1 \times 9}, \kappa = 10, \delta = 10$.
    \end{enumerate}
     The following scenarios in the circle are considered
    \begin{enumerate}
        \item $r=2, d=1, \bmu_1 = (\sqrt{1/2})_{1 \times 4}, \kappa_1 = 0, \kappa = 0$ (independence case);
        \item $r=2, d=1, \bmu_1 = (\sqrt{1/2})_{1 \times 4}, \kappa_1 = 0, \kappa = 1$;
        \item $r=2, d=1, \bmu_1 = (\sqrt{1/2})_{1 \times 4}, \kappa_1 = 0, \kappa = 2$;
        \item $r=2, d=1, \bmu_1 = (\sqrt{1/2})_{1 \times 4}, \kappa_1 = 0, \kappa = 5$;
        \item $r=2, d=1, \bmu_1 = (\sqrt{1/2})_{1 \times 4}, \kappa_1 = 0, \kappa = 10$.
    \end{enumerate}
    \item \textbf{Markovian process:} we simulate $\bX^{(1)} \sim \mathrm{vMF}_d (\bmu_1, \kappa_1)$, and, then, we simulate $\bX^{(j)} | (\bX^{(1)} = \bx^{(1)} , \dots, \bX^{(j-1)} = \bx^{(j-1)}) \sim \mathrm{vMF}_d (\bx^{(j-1)}, \delta)$ for $j = 2, \dots, r$. This produces the density
    \begin{align*}
        f(\bx) &= c^\mathrm{vMF}_{d}(\kappa_1) e^{\kappa_1\bx^{(1) \top}\bmu_1} \prod_{j=2}^{r} c^\mathrm{vMF}_{d}(\delta) e^{\delta\bx^{(j) \top}\bx^{(j-1)}} \\
        &= c^\mathrm{vMF}_{d}(\kappa_1) c^\mathrm{vMF}_{d}(\kappa)^{r-1} e^{\kappa_1\bx^{(1) \top}\bmu_1 + \sum_{j=2}^{r} \delta\bx^{(j) \top}\bx^{(j-1)}}.
\end{align*}

    The following scenarios will be considered.
    \begin{enumerate}
        \item $r=10, d=1, \bmu_1 = (\sqrt{1/2})_{1 \times 20}, \kappa_1 = 1, \delta = 0$ (independence case);
        \item $r=10, d=1, \bmu_1 = (\sqrt{1/2})_{1 \times 20}, \kappa_1 = 1, \delta = 0.1$;
        \item $r=10, d=1, \bmu_1 = (\sqrt{1/2})_{1 \times 20}, \kappa_1 = 1, \delta = 0.5$;
        \item $r=10, d=1, \bmu_1 = (\sqrt{1/2})_{1 \times 20}, \kappa_1 = 1, \delta = 1$;
        \item $r=10, d=1, \bmu_1 = (\sqrt{1/2})_{1 \times 20}, \kappa_1 = 1, \delta = 10$.
    \end{enumerate}
\end{enumerate}

Scenario i is constructed based on the fact that conditional independence is a simple way to induce dependence between random vectors. Note that, as a particular case, when the distribution of the latent random vector $\bmu$ is degenerate (which could be interpreted as $\kappa_1 \to \infty$ or $\delta = 0$), the random vectors $\bX^{(1)}, \dots, \bX^{(r)}$ would be unconditionally iid.

Scenario ii produces a sequence of random vectors, which of them supported in a hypersphere, where the location parameter of one component $\bX^{(j)}$ given the previous ones is given by the last component $\bX^{(j-1)}$. As in the first scenario, the parameter $\delta$ represents the amount of deviation from independence.

Each scenario is simulated $100$ times for $n = 10, 50, 100$, using $B=100$ replicas in each p-value computation. Notice that scenario (a) represents the null hypothesis because, under $\delta = 0$, the components $\bX^{(1)}, \dots, \bX^{(r)}$ are independent. The adopted decision rule is to reject the null hypothesis when the p-value is less or equal to $\alpha = 0.05$. The bandwidths used in the simulation were chosen by likelihood cross-validation (see \cite{Hall1987}).

The results of the simulation are summarized in Table \ref{tab1}. As expected, the power increased as a function of $n$ - except for $\delta = 0$, which is the independence case. In such a case, the empirical power is close to 0.05. Furthermore, the rejection rate has increased as a function of $\delta$, which is reasonable, since greater values of $\delta$ represent a stronger deviation from independence.

\begin{table}[h!]
\centering
\caption{Empirical power for each simulated case (scenario i), adopting $\alpha = 0.05$.}
\label{tab:results}
\begin{tabular}{lccc}
\toprule
    & $n=10$ & $n=50$ & $n=100$ \\
\midrule
(a) & 0.06 & 0.03 & 0.05 \\
(b) & 0.25 & 0.85 & 0.99 \\
(c) & 0.82 & 1.00 & 1.00 \\
(d) & 0.95 & 1.00 & 1.00 \\
(e) & 0.98 & 1.00 & 1.00 \\
\bottomrule
\end{tabular}
\label{tab1}
\end{table}

\begin{table}[ht]
\centering
\caption{Empirical power for each simulated case (scenario i, $r=2, d = 1$), adopting $\alpha = 0.05$.}
\label{tab:results}
\begin{tabular}{lccc}
\toprule
      & $n=10$ & $n=50$ & $n=100$ \\
\midrule
(a)   & 0.04 & 0.07 & 0.09 \\
(b)   & 0.07 & 0.08 & 0.04 \\
(c)   & 0.04 & 0.12 & 0.06 \\
(d)   & 0.08 & 0.21 & 0.52 \\
(e)   & 0.84 & 1.00 & 1.00 \\
\bottomrule
\end{tabular}
\end{table}

\begin{table}[H]
\centering
\caption{Empirical power for each simulated case (scenario i, $r=2, d = 1$), adopting $\alpha = 0.05$ ($\kappa = 0,1,2,5,10$).}
\label{tab:results}
\begin{tabular}{lccc}
\toprule
      & $n=10$ & $n=50$ & $n=100$ \\
\midrule
(a) & 0.04 & 0.07 & 0.09  \\
(b) & 0.08 & 0.21 & 0.52  \\
(c) & 0.3  & 0.94 & 1     \\
(d) & 0.75 & 1    & 1     \\
(e) & 0.84 & 1    & 1    \\
\bottomrule
\end{tabular}
\end{table}

\begin{figure}[H]
\centering
        \includegraphics[totalheight=8cm]{Rplot.png}
    \caption{$\kappa = 0$ (independence case).}
\end{figure}

\begin{figure}[H]
\centering
        \includegraphics[totalheight=8cm]{Rplot01.png}
    \caption{$\kappa = 1$.}
\end{figure}

\begin{figure}[H]
\centering
        \includegraphics[totalheight=8cm]{Rplot02.png}
    \caption{$\kappa = 2$.}
\end{figure}

\begin{figure}[H]
\centering
        \includegraphics[totalheight=8cm]{Rplot03.png}
    \caption{$\kappa = 5$.}
\end{figure}

\begin{figure}[H]
\centering
        \includegraphics[totalheight=8cm]{Rplot04.png}
    \caption{$\kappa = 10$.}
\end{figure}

\begin{table}[ht]
\centering
\caption{Empirical power for $d=\cdots$, $r=\cdots$ (nominal level $0.05$).}
\begin{tabular}{ccccccc}
\toprule
 &  & \multicolumn{5}{c}{$\kappa$} \\
\cmidrule(lr){3-7}
$n$ & $h$ 
& $\kappa_1$ & $\kappa_2$ & $\kappa_3$ & $\kappa_4$ & $\kappa_5$ \\
\midrule

% n = n1
$n_1$ & $h_1$ &  &  &  &  &  \\
      & $h_2$ &  &  &  &  &  \\
      & $h_3$ &  &  &  &  &  \\
      & $h_4$ &  &  &  &  &  \\
      & $h_5$ &  &  &  &  &  \\
\midrule

% n = n2
$n_2$ & $h_1$ &  &  &  &  &  \\
      & $h_2$ &  &  &  &  &  \\
      & $h_3$ &  &  &  &  &  \\
      & $h_4$ &  &  &  &  &  \\
      & $h_5$ &  &  &  &  &  \\
\midrule

% n = n3
$n_3$ & $h_1$ &  &  &  &  &  \\
      & $h_2$ &  &  &  &  &  \\
      & $h_3$ &  &  &  &  &  \\
      & $h_4$ &  &  &  &  &  \\
      & $h_5$ &  &  &  &  &  \\
\midrule

% n = n4
$n_4$ & $h_1$ &  &  &  &  &  \\
      & $h_2$ &  &  &  &  &  \\
      & $h_3$ &  &  &  &  &  \\
      & $h_4$ &  &  &  &  &  \\
      & $h_5$ &  &  &  &  &  \\
\midrule

% n = n5
$n_5$ & $h_1$ &  &  &  &  &  \\
      & $h_2$ &  &  &  &  &  \\
      & $h_3$ &  &  &  &  &  \\
      & $h_4$ &  &  &  &  &  \\
      & $h_5$ &  &  &  &  &  \\

\bottomrule
\end{tabular}
\end{table}

% latex table generated in R 4.5.2 by xtable 1.8-8 package
% Mon Feb 23 16:35:39 2026
\begin{table}[ht]
\centering
\begin{tabular}{lrrrr}
  \hline
scenario & kappa\_0 & kappa\_0.333333333333333 & kappa\_1 & kappa\_3 \\ 
  \hline
n=10, c=0,& 0.03 & 0.05 & 0.09 & 0.66 \\ 
  n=50, c=0,& 0.05 & 0.06 & 0.34 & 1.00 \\ 
  n=100, c=0,& 0.05 & 0.06 & 0.62 & 1.00 \\ 
  n=10, c=1,& 0.04 & 0.05 & 0.09 & 0.67 \\ 
  n=50, c=1,& 0.05 & 0.06 & 0.29 & 1.00 \\ 
  n=100, c=1,& 0.04 & 0.06 & 0.60 & 1.00 \\ 
  n=10, c=2,& 0.04 & 0.04 & 0.07 & 0.69 \\ 
  n=50, c=2,& 0.05 & 0.05 & 0.28 & 1.00 \\ 
  n=100, c=2,& 0.04 & 0.06 & 0.59 & 1.00 \\ 
  n=10, c=3,& 0.06 & 0.05 & 0.09 & 0.70 \\ 
  n=50, c=3,& 0.05 & 0.06 & 0.30 & 1.00 \\ 
  n=100, c=3,& 0.04 & 0.05 & 0.59 & 1.00 \\ 
   \hline
\end{tabular}
\caption{Empirical power for alpha=0.05 across different scenarios and kappa values.} 
\end{table}

\begin{table}[ht]
\centering
\begin{tabular}{cc S S S S}
\toprule
$n$ & $c$ & {$\kappa=0$} & {$\kappa=\tfrac{1}{3}$} & {$\kappa=1$} & {$\kappa=3$} \\
\midrule
10  & 0 & 0.03 & 0.05 & 0.09 & 0.66 \\
50  & 0 & 0.05 & 0.06 & 0.34 & 1.00 \\
100 & 0 & 0.05 & 0.06 & 0.62 & 1.00 \\
\addlinespace
10  & 1 & 0.04 & 0.05 & 0.09 & 0.67 \\
50  & 1 & 0.05 & 0.06 & 0.29 & 1.00 \\
100 & 1 & 0.04 & 0.06 & 0.60 & 1.00 \\
\addlinespace
10  & 2 & 0.04 & 0.04 & 0.07 & 0.69 \\
50  & 2 & 0.05 & 0.05 & 0.28 & 1.00 \\
100 & 2 & 0.04 & 0.06 & 0.59 & 1.00 \\
\addlinespace
10  & 3 & 0.06 & 0.05 & 0.09 & 0.70 \\
50  & 3 & 0.05 & 0.06 & 0.30 & 1.00 \\
100 & 3 & 0.04 & 0.05 & 0.59 & 1.00 \\
\bottomrule
\end{tabular}
\caption{Empirical power ($\alpha=0.05$) across sample sizes, tuning parameter $c$, and signal strength $\kappa$.}
\end{table}

\section{New table}

\begin{table}[!h]
\centering
\caption{Empirical power at $\alpha = 0.05$ for all simulated scenarios.}
\centering
\resizebox{\ifdim\width>\linewidth\linewidth\else\width\fi}{!}{
\fontsize{10}{12}\selectfont
\begin{tabular}[t]{cclcccc}
\toprule
\multicolumn{3}{c}{ } & \multicolumn{4}{c}{Dependence strength} \\
\cmidrule(l{3pt}r{3pt}){4-7}
\textbf{d} & \textbf{r} & \textbf{c} & \textbf{$\kappa = 0$} & \textbf{$\kappa = 1$} & \textbf{$\kappa = 2$} & \textbf{$\kappa = 4$}\\
\midrule
 &  & $2^{-1}h_{\mathrm{ideal}}$ & 0.060 & 0.375 & 1.000 & 1.000\\

 &  & $h_{\mathrm{ideal}}$ & 0.060 & 0.430 & 1.000 & 1.000\\

 & \multirow{-3}{*}{\centering\arraybackslash 2} & $2h_{\mathrm{ideal}}$ & 0.060 & 0.430 & 1.000 & 1.000\\

 &  & $2^{-1}h_{\mathrm{ideal}}$ & 0.035 & 0.665 & 1.000 & 1.000\\

 &  & $h_{\mathrm{ideal}}$ & 0.030 & 0.715 & 1.000 & 1.000\\

 & \multirow{-3}{*}{\centering\arraybackslash 3} & $2h_{\mathrm{ideal}}$ & 0.030 & 0.695 & 1.000 & 1.000\\

 &  & $2^{-1}h_{\mathrm{ideal}}$ & 0.045 & 0.315 & 1.000 & 1.000\\

 &  & $h_{\mathrm{ideal}}$ & 0.060 & 1.000 & 1.000 & 1.000\\

\multirow{-9}{*}[1\dimexpr\aboverulesep+\belowrulesep+\cmidrulewidth]{\centering\arraybackslash 1} & \multirow{-3}{*}{\centering\arraybackslash 10} & $2h_{\mathrm{ideal}}$ & 0.080 & 1.000 & 1.000 & 1.000\\
\cmidrule{1-7}
 &  & $2^{-1}h_{\mathrm{ideal}}$ & 0.065 & 0.100 & 0.840 & 1.000\\

 &  & $h_{\mathrm{ideal}}$ & 0.050 & 0.135 & 0.935 & 1.000\\

 & \multirow{-3}{*}{\centering\arraybackslash 2} & $2h_{\mathrm{ideal}}$ & 0.045 & 0.115 & 0.930 & 1.000\\

 &  & $2^{-1}h_{\mathrm{ideal}}$ & 0.025 & 0.195 & 0.990 & 1.000\\

 &  & $h_{\mathrm{ideal}}$ & 0.040 & 0.225 & 1.000 & 1.000\\

 & \multirow{-3}{*}{\centering\arraybackslash 3} & $2h_{\mathrm{ideal}}$ & 0.030 & 0.230 & 1.000 & 1.000\\

 &  & $2^{-1}h_{\mathrm{ideal}}$ & 0.055 & 0.060 & 0.075 & 0.080\\

 &  & $h_{\mathrm{ideal}}$ & 0.060 & 0.710 & 1.000 & 1.000\\

\multirow{-9}{*}[1\dimexpr\aboverulesep+\belowrulesep+\cmidrulewidth]{\centering\arraybackslash 2} & \multirow{-3}{*}{\centering\arraybackslash 10} & $2h_{\mathrm{ideal}}$ & 0.070 & 0.800 & 1.000 & 1.000\\
\cmidrule{1-7}
 &  & $2^{-1}h_{\mathrm{ideal}}$ & 0.030 & 0.055 & 0.045 & 0.050\\

 &  & $h_{\mathrm{ideal}}$ & 0.020 & 0.055 & 0.060 & 0.135\\

 & \multirow{-3}{*}{\centering\arraybackslash 2} & $2h_{\mathrm{ideal}}$ & 0.035 & 0.045 & 0.070 & 0.180\\

 &  & $2^{-1}h_{\mathrm{ideal}}$ & 0.050 & 0.055 & 0.040 & 0.065\\

 &  & $h_{\mathrm{ideal}}$ & 0.065 & 0.040 & 0.025 & 0.155\\

 & \multirow{-3}{*}{\centering\arraybackslash 3} & $2h_{\mathrm{ideal}}$ & 0.050 & 0.050 & 0.065 & 0.310\\

 &  & $2^{-1}h_{\mathrm{ideal}}$ & 0.035 & 0.075 & 0.090 & 0.055\\

 &  & $h_{\mathrm{ideal}}$ & 0.060 & 0.045 & 0.065 & 0.050\\

\multirow{-9}{*}[1\dimexpr\aboverulesep+\belowrulesep+\cmidrulewidth]{\centering\arraybackslash 10} & \multirow{-3}{*}{\centering\arraybackslash 10} & $2h_{\mathrm{ideal}}$ & 0.025 & 0.075 & 0.085 & 0.390\\
\bottomrule
\end{tabular}}
\end{table}

\section{New table (Markov)}

\begin{table}[!h]
\centering
\caption{Empirical power at $\alpha = 0.05$ for all simulated scenarios (Markov model).}
\centering
\resizebox{\ifdim\width>\linewidth\linewidth\else\width\fi}{!}{
\fontsize{10}{12}\selectfont
\begin{tabular}[t]{cclcccc}
\toprule
\multicolumn{3}{c}{ } & \multicolumn{4}{c}{Dependence strength} \\
\cmidrule(l{3pt}r{3pt}){4-7}
\textbf{d} & \textbf{r} & \textbf{c} & \textbf{$\kappa = 0$} & \textbf{$\kappa = 1$} & \textbf{$\kappa = 2$} & \textbf{$\kappa = 4$}\\
\midrule
 &  & $2^{-1}h_{\mathrm{ideal}}$ & 0.060 & 1.000 & 1.000 & 1.000\\

 &  & $h_{\mathrm{ideal}}$ & 0.040 & 1.000 & 1.000 & 1.000\\

 & \multirow{-3}{*}{\centering\arraybackslash 2} & $2h_{\mathrm{ideal}}$ & 0.060 & 1.000 & 1.000 & 1.000\\

 &  & $2^{-1}h_{\mathrm{ideal}}$ & 0.085 & 1.000 & 1.000 & 1.000\\

 &  & $h_{\mathrm{ideal}}$ & 0.065 & 1.000 & 1.000 & 1.000\\

 & \multirow{-3}{*}{\centering\arraybackslash 3} & $2h_{\mathrm{ideal}}$ & 0.065 & 1.000 & 1.000 & 1.000\\

 &  & $2^{-1}h_{\mathrm{ideal}}$ & 0.060 & 0.985 & 1.000 & 0.565\\

 &  & $h_{\mathrm{ideal}}$ & 0.040 & 1.000 & 1.000 & 1.000\\

\multirow{-9}{*}[1\dimexpr\aboverulesep+\belowrulesep+\cmidrulewidth]{\centering\arraybackslash 1} & \multirow{-3}{*}{\centering\arraybackslash 10} & $2h_{\mathrm{ideal}}$ & 0.060 & 1.000 & 1.000 & 1.000\\
\cmidrule{1-7}
 &  & $2^{-1}h_{\mathrm{ideal}}$ & 0.025 & 0.890 & 1.000 & 1.000\\

 &  & $h_{\mathrm{ideal}}$ & 0.050 & 0.975 & 1.000 & 1.000\\

 & \multirow{-3}{*}{\centering\arraybackslash 2} & $2h_{\mathrm{ideal}}$ & 0.045 & 0.955 & 1.000 & 1.000\\

 &  & $2^{-1}h_{\mathrm{ideal}}$ & 0.040 & 0.905 & 1.000 & 1.000\\

 &  & $h_{\mathrm{ideal}}$ & 0.070 & 0.985 & 1.000 & 1.000\\

 & \multirow{-3}{*}{\centering\arraybackslash 3} & $2h_{\mathrm{ideal}}$ & 0.060 & 0.990 & 1.000 & 1.000\\

 &  & $2^{-1}h_{\mathrm{ideal}}$ & 0.035 & 0.060 & 0.070 & 0.050\\

 &  & $h_{\mathrm{ideal}}$ & 0.035 & 0.960 & 1.000 & 1.000\\

\multirow{-9}{*}[1\dimexpr\aboverulesep+\belowrulesep+\cmidrulewidth]{\centering\arraybackslash 2} & \multirow{-3}{*}{\centering\arraybackslash 10} & $2h_{\mathrm{ideal}}$ & 0.030 & 0.980 & 1.000 & 1.000\\
\cmidrule{1-7}
 &  & $2^{-1}h_{\mathrm{ideal}}$ & 0.040 & 0.055 & 0.040 & 0.055\\

 &  & $h_{\mathrm{ideal}}$ & 0.050 & 0.115 & 0.360 & 1.000\\

 & \multirow{-3}{*}{\centering\arraybackslash 2} & $2h_{\mathrm{ideal}}$ & 0.045 & 0.130 & 0.540 & 1.000\\

 &  & $2^{-1}h_{\mathrm{ideal}}$ & 0.035 & 0.070 & 0.035 & 0.085\\

 &  & $h_{\mathrm{ideal}}$ & 0.060 & 0.065 & 0.160 & 0.515\\

 & \multirow{-3}{*}{\centering\arraybackslash 3} & $2h_{\mathrm{ideal}}$ & 0.045 & 0.100 & 0.470 & 1.000\\

 &  & $2^{-1}h_{\mathrm{ideal}}$ & 0.030 & 0.050 & 0.080 & 0.050\\

 &  & $h_{\mathrm{ideal}}$ & 0.040 & 0.030 & 0.025 & 0.090\\

\multirow{-9}{*}[1\dimexpr\aboverulesep+\belowrulesep+\cmidrulewidth]{\centering\arraybackslash 10} & \multirow{-3}{*}{\centering\arraybackslash 10} & $2h_{\mathrm{ideal}}$ & 0.050 & 0.060 & 0.215 & 0.905\\
\bottomrule
\end{tabular}}
\end{table}

\section{Expectation of G}

\paragraph{Asymptotic behavior of $E(\widetilde G_n^2(X_1,X_2))$.}
Recall that
\[
\widetilde H_n(x,X_1):=\int_{S^d}\bar L_h(t,x)\bar L_h(t,X_1)\,dt,
\]
and
\[
\widetilde G_n(x,y):=E\!\left[\widetilde H_n(x,X_1)\widetilde H_n(y,X_1)\right].
\]
Hence,
\[
E\!\left[\widetilde G_n^2(X_1,X_2)\right]
=
\int_{S^d}\int_{S^d}\widetilde G_n(x,y)^2 f(x)f(y)\,dx\,dy.
\]

Expanding the centered kernels, $\bar L_h=L_h-E(L_h)$, the quantity
$E(\widetilde G_n^2(X_1,X_2))$ is split into $16$ summands. We show that the first one is of
order $O(\rho(h)^{-1})$ and that the remaining terms are of smaller order. The first summand is
\[
C_1:=
\int_{S^d}\int_{S^d}
\left\{
\int_{S^d}\int_{S^d}\int_{S^d}
L_h(t,x)L_h(t,u)L_h(s,y)L_h(s,u)f(u)\,dt\,ds\,du
\right\}^2
f(x)f(y)\,dx\,dy.
\]

Define the overlap kernel
\[
H_h^*(x,u):=\int_{S^d}L_h(t,x)L_h(t,u)\,dt.
\]
Then
\[
C_1=
\int_{S^d}\int_{S^d}
\left\{
\int_{S^d}H_h^*(x,u)H_h^*(y,u)f(u)\,du
\right\}^2
f(x)f(y)\,dx\,dy.
\]

Under the null hypothesis, $f(x)=\prod_{j=1}^r f_j(x^{(j)})$, and, because $L_h$ is a product kernel,
\[
H_h^*(x,u)=\prod_{j=1}^r \varphi_{j,h_j}(x^{(j)},u^{(j)}),
\]
where
\[
\varphi_{j,h_j}(a,b):=
\int_{S^{d_j}}L_{h_j}(z,a)L_{h_j}(z,b)\,dz.
\]
Therefore,
\[
\int_{S^d}H_h^*(x,u)H_h^*(y,u)f(u)\,du
=
\prod_{j=1}^r
\left\{
\int_{S^{d_j}}
\varphi_{j,h_j}(x^{(j)},u^{(j)})
\varphi_{j,h_j}(y^{(j)},u^{(j)})
f_j(u^{(j)})\,du^{(j)}
\right\},
\]
and, consequently,
\[
C_1=\prod_{j=1}^r C_{1,j},
\]
with
\[
C_{1,j}:=
\int_{S^{d_j}}\int_{S^{d_j}}
\left\{
\int_{S^{d_j}}
\varphi_{j,h_j}(x,u)\varphi_{j,h_j}(y,u)f_j(u)\,du
\right\}^2
f_j(x)f_j(y)\,dx\,dy.
\]

Now we study $\varphi_{j,h_j}$. By definition,
\[
\varphi_{j,h_j}(a,b)
=
c_{d_j,L}(h_j)^2
\int_{S^{d_j}}
L\!\left(\frac{1-z^\top a}{h_j^2}\right)
L\!\left(\frac{1-z^\top b}{h_j^2}\right)\,dz.
\]
Applying the tangent-normal change of variables in Lemma B.4, exactly as in the derivation of the
dominant term $A_1$ in the proof of $E(\widetilde H_n^2(X_1,X_2))$, it follows that there exists a
bounded kernel $\psi_j:\mathbb R_+\to\mathbb R_+$ such that
\[
\varphi_{j,h_j}(a,b)
=
h_j^{-d_j}\psi_j\!\left(\frac{1-a^\top b}{h_j^2}\right)(1+o(1)),
\qquad h_j\to0,
\]
uniformly on $S^{d_j}\times S^{d_j}$. In particular, $\varphi_{j,h_j}$ is itself a directional kernel
with normalizing order $h_j^{-d_j}$.

Therefore, applying again the same argument used for the dominant term $A_1$ in the proof of
$E(\widetilde H_n^2(X_1,X_2))$, but now with $\varphi_{j,h_j}$ in place of $L_{h_j}$, yields
\[
C_{1,j}=O(h_j^{-d_j}).
\]
Hence,
\[
C_1
=
\prod_{j=1}^r O(h_j^{-d_j})
=
O\!\left(\prod_{j=1}^r h_j^{-d_j}\right)
=
O(\rho(h)^{-1}).
\]

Consider now any of the remaining $15$ summands in the expansion of
$E(\widetilde G_n^2(X_1,X_2))$. Each of them contains at least one factor of the form
$E\{L_h(\cdot,X)\}$ instead of a genuine overlap term. By Remark 2.1 and Lemma B.3,
\[
E\{L_h(x,X)\}=f(x)+O\!\left(\sum_{j=1}^r h_j^2\right)=O(1)
\]
uniformly in $x$, and therefore each such term loses at least one overlap-kernel contribution.
Repeating the same iterative argument as above shows that none of these remaining terms exceeds
the order of $C_1$; in fact, they are all $O(1)$ or $o(\rho(h)^{-1})$.

Thus,
\[
E\!\left[\widetilde G_n^2(X_1,X_2)\right]=O(\rho(h)^{-1}).
\]

\subsection{New proof}

\paragraph{Asymptotic behavior of $E(\widetilde G_n^2(X_1,X_2))$.}
Recall that
\[
\widetilde G_n(x,y):=E\!\left[\widetilde H_n(x,X_1)\widetilde H_n(y,X_1)\right],
\]
with
\[
\widetilde H_n(x,u):=\int_{S^d}\overline L_h(z,x)\overline L_h(z,u)\,dz,
\qquad
\overline L_h(z,x):=L_h(z,x)-E[L_h(z,X_1)].
\]
Therefore,
\[
E\!\left[\widetilde G_n^2(X_1,X_2)\right]
=
\int_{S^d}\int_{S^d}\widetilde G_n(x,y)^2 f(x)f(y)\,dx\,dy.
\]

Define
\[
m_h(x):=E[L_h(x,X_1)],
\qquad
H_n^*(x,u):=\int_{S^d}L_h(z,x)L_h(z,u)\,dz,
\]
and also
\[
B_h(x):=\int_{S^d}L_h(z,x)m_h(z)\,dz,
\qquad
C_h:=\int_{S^d}m_h(z)^2\,dz.
\]
Then
\[
\widetilde H_n(x,u)=H_n^*(x,u)-B_h(x)-B_h(u)+C_h.
\]
Hence,
\[
\widetilde G_n(x,y)
=
A_h(x,y)+R_h(x,y),
\]
where
\[
A_h(x,y):=E\!\left[H_n^*(x,X_1)H_n^*(y,X_1)\right]
=
\int_{S^d}H_n^*(x,u)H_n^*(y,u)f(u)\,du,
\]
and $R_h(x,y)$ collects all the remaining terms.

By Remark 2.1,
\[
m_h(x)=f(x)+O\!\left(\sum_{j=1}^r h_j^2\right)
\]
uniformly in $x\in S^d$. Since $L_h(\cdot,\cdot)$ is a probability kernel, it follows that
$B_h(x)=O(1)$ uniformly in $x$, and also $C_h=O(1)$. Moreover,
\[
E\!\left[H_n^*(x,X_1)\right]
=
\int_{S^d}L_h(z,x)m_h(z)\,dz
=
B_h(x)
=
O(1)
\]
uniformly in $x$. Therefore,
\[
R_h(x,y)=O(1)
\]
uniformly in $(x,y)\in S^d\times S^d$, and consequently
\[
\int_{S^d}\int_{S^d}R_h(x,y)^2 f(x)f(y)\,dx\,dy=O(1).
\]
Thus,
\[
E\!\left[\widetilde G_n^2(X_1,X_2)\right]
=
\int\!\!\int \big(A_h(x,y)+R_h(x,y)\big)^2 f(x)f(y)\,dx\,dy
\]
\[
\le
2\int\!\!\int A_h(x,y)^2 f(x)f(y)\,dx\,dy
+
2\int\!\!\int R_h(x,y)^2 f(x)f(y)\,dx\,dy
=
2A_1+O(1),
\]
where
\[
A_1:=
\int_{S^d}\int_{S^d}
\left(
\int_{S^d}H_n^*(x,u)H_n^*(y,u)f(u)\,du
\right)^2
f(x)f(y)\,dx\,dy.
\]

It remains to study $A_1$. Since $L_h$ is a product kernel, then
\[
H_n^*(x,u)
=
\prod_{j=1}^r
\varphi_{j,h_j}\!\left(x^{(j)},u^{(j)}\right),
\]
where
\[
\varphi_{j,h_j}(a,b):=
\int_{S^{d_j}}L_{h_j}(z,a)L_{h_j}(z,b)\,dz.
\]
Under the null hypothesis,
\[
f(u)=\prod_{j=1}^r f_j(u^{(j)}),
\]
and therefore
\[
A_1
=
\prod_{j=1}^r A_{1,j},
\]
with
\[
A_{1,j}:=
\int_{S^{d_j}}\int_{S^{d_j}}
\left(
\int_{S^{d_j}}
\varphi_{j,h_j}(x,u)\varphi_{j,h_j}(y,u)f_j(u)\,du
\right)^2
f_j(x)f_j(y)\,dx\,dy.
\]

Now, for each $j$, the function $\varphi_{j,h_j}$ is a kernel-overlap term of the same type as the
quantity studied in the previous subsection for $H_n^*(X_1,X_2)$. Repeating the same tangent-normal
change-of-variables argument componentwise, and using that
\[
c_{d_j,L}(h_j)\asymp h_j^{-d_j},
\]
one obtains
\[
A_{1,j}=O(h_j^{-d_j}).
\]
Hence,
\[
A_1
=
\prod_{j=1}^r O(h_j^{-d_j})
=
O\!\left(\prod_{j=1}^r h_j^{-d_j}\right)
=
O(\rho(h)^{-1}).
\]

Therefore,
\[
E\!\left[\widetilde G_n^2(X_1,X_2)\right]
\le
2A_1+O(1)
=
O(\rho(h)^{-1}),
\]
which proves the claim.

\subsection{Other new proof}

\paragraph{Asymptotic behavior of $\mathsf{E}\lrc{G_n^2(\bX_1,\bX_2)}$.}
Recall that
\[
\bar H_n(\bx,\bu) \defin \int_{\Sr} \bar L_{\bh}(\bz,\bx)\bar L_{\bh}(\bz,\bu)\,\rd \bz,
\qquad
\bar L_{\bh}(\bz,\bx) \defin L_{\bh}(\bz,\bx)-\mathsf{E}\lrc{L_{\bh}(\bz,\bX_1)},
\]
and
\[
G_n(\bx,\by) \defin \mathsf{E}\lrc{\bar H_n(\bx,\bX_1)\bar H_n(\by,\bX_1)}.
\]
Therefore,
\[
\mathsf{E}\lrc{G_n^2(\bX_1,\bX_2)}
=
\int_{\Sr}\int_{\Sr} G_n(\bx,\by)^2 f(\bx)f(\by)\,\rd \bx\,\rd \by.
\]

Define
\[
m_n(\bz) \defin \mathsf{E}\lrc{L_{\bh}(\bz,\bX_1)}
=
\int_{\Sr} L_{\bh}(\bz,\bv)f(\bv)\,\rd \bv,
\]
\[
H_n^\ast(\bx,\bu) \defin \int_{\Sr} L_{\bh}(\bz,\bx)L_{\bh}(\bz,\bu)\,\rd \bz,
\qquad
B_n(\bx) \defin \int_{\Sr} L_{\bh}(\bz,\bx)m_n(\bz)\,\rd \bz,
\qquad
C_n \defin \int_{\Sr} m_n(\bz)^2\,\rd \bz.
\]
Then
\[
\bar H_n(\bx,\bu)=H_n^\ast(\bx,\bu)-B_n(\bx)-B_n(\bu)+C_n.
\]
Hence,
\[
G_n(\bx,\by)=A_n(\bx,\by)+R_n(\bx,\by),
\]
where
\[
A_n(\bx,\by)\defin \int_{\Sr} H_n^\ast(\bx,\bu)H_n^\ast(\by,\bu)f(\bu)\,\rd \bu,
\]
and \(R_n(\bx,\by)\) collects the remaining terms.

Since \(f\) is continuous on the compact set \(\Sr\), it is bounded. Thus,
\[
0 \le m_n(\bz)=\int_{\Sr}L_{\bh}(\bz,\bv)f(\bv)\,\rd \bv
\le \norm{f}_\infty \int_{\Sr}L_{\bh}(\bz,\bv)\,\rd \bv
= \norm{f}_\infty,
\]
so \(m_n(\bz)=O(1)\) uniformly in \(\bz\in\Sr\). Consequently,
\[
B_n(\bx)=\int_{\Sr}L_{\bh}(\bz,\bx)m_n(\bz)\,\rd \bz=O(1)
\]
uniformly in \(\bx\), and also
\[
C_n=\int_{\Sr}m_n(\bz)^2\,\rd \bz=O(1).
\]
Moreover,
\[
\int_{\Sr} H_n^\ast(\bx,\bu)f(\bu)\,\rd \bu
=
\int_{\Sr}\int_{\Sr} L_{\bh}(\bz,\bx)L_{\bh}(\bz,\bu)f(\bu)\,\rd \bz\,\rd \bu
=
\int_{\Sr}L_{\bh}(\bz,\bx)m_n(\bz)\,\rd \bz
=
B_n(\bx),
\]
hence
\[
\int_{\Sr} H_n^\ast(\bx,\bu)f(\bu)\,\rd \bu=O(1)
\]
uniformly in \(\bx\). Therefore, every term entering \(R_n(\bx,\by)\) is uniformly \(O(1)\), which
implies
\[
\int_{\Sr}\int_{\Sr}R_n(\bx,\by)^2f(\bx)f(\by)\,\rd \bx\,\rd \by=O(1).
\]

By \((a+b)^2\leq 2a^2+2b^2\),
\[
\mathsf{E}\lrc{G_n^2(\bX_1,\bX_2)}
\leq
2\int_{\Sr}\int_{\Sr}A_n(\bx,\by)^2f(\bx)f(\by)\,\rd \bx\,\rd \by
+
2\int_{\Sr}\int_{\Sr}R_n(\bx,\by)^2f(\bx)f(\by)\,\rd \bx\,\rd \by.
\]
Hence,
\[
\mathsf{E}\lrc{G_n^2(\bX_1,\bX_2)} \leq 2A_{n,1}+O(1),
\]
where
\[
A_{n,1}
\defin
\int_{\Sr}\int_{\Sr}
\lrp{
\int_{\Sr}H_n^\ast(\bx,\bu)H_n^\ast(\by,\bu)f(\bu)\,\rd \bu
}^2
f(\bx)f(\by)\,\rd \bx\,\rd \by.
\]

Using again that \(f\) is bounded,
\[
\lrp{
\int_{\Sr}H_n^\ast(\bx,\bu)H_n^\ast(\by,\bu)f(\bu)\,\rd \bu
}^2
\leq
\norm{f}_\infty^2
\lrp{
\int_{\Sr}H_n^\ast(\bx,\bu)H_n^\ast(\by,\bu)\,\rd \bu
}^2.
\]
Therefore,
\[
A_{n,1}
\leq
\norm{f}_\infty^4
\int_{\Sr}\int_{\Sr}Q_n(\bx,\by)^2\,\rd \bx\,\rd \by,
\]
where
\[
Q_n(\bx,\by)\defin \int_{\Sr}H_n^\ast(\bx,\bu)H_n^\ast(\by,\bu)\,\rd \bu.
\]

Now we exploit the product structure of the kernel. Since
\[
L_{\bh}(\bx,\bu)=\prod_{j=1}^r L_{h_j}\lrp{\bx^{(j)},\bu^{(j)}},
\]
it follows that
\[
H_n^\ast(\bx,\bu)=\prod_{j=1}^r \varphi_{j,h_j}\lrp{\bx^{(j)},\bu^{(j)}},
\]
where
\[
\varphi_{j,h_j}(\ba,\bb)
\defin
\int_{\Sj}L_{h_j}(\bz,\ba)L_{h_j}(\bz,\bb)\,\rd \bz.
\]
Hence,
\[
Q_n(\bx,\by)=\prod_{j=1}^r q_{j,h_j}\lrp{\bx^{(j)},\by^{(j)}},
\]
with
\[
q_{j,h_j}(\ba,\bb)
\defin
\int_{\Sj}\varphi_{j,h_j}(\ba,\bu)\varphi_{j,h_j}(\bb,\bu)\,\rd \bu.
\]

Fix \(j\in\lrb{1,\ldots,r}\). First,
\[
\int_{\Sj}\varphi_{j,h_j}(\ba,\bu)\,\rd \bu
=
\int_{\Sj}\int_{\Sj}L_{h_j}(\bz,\ba)L_{h_j}(\bz,\bu)\,\rd \bz\,\rd \bu
=
\int_{\Sj}L_{h_j}(\bz,\ba)\,\rd \bz
=
1.
\]
Also, since \(L\) is bounded and \(c_{d_j,L}(h_j)\asymp h_j^{-d_j}\),
\[
\sup_{\ba,\bu}\varphi_{j,h_j}(\ba,\bu)
\leq
\sup_{\bz,\ba}L_{h_j}(\bz,\ba)
=
O(h_j^{-d_j}).
\]
Therefore,
\[
0\leq q_{j,h_j}(\ba,\bb)
\leq
\sup_{\bu} \varphi_{j,h_j}(\ba,\bu)\int_{\Sj}\varphi_{j,h_j}(\bb,\bu)\,\rd \bu
=
O(h_j^{-d_j}),
\]
uniformly in \((\ba,\bb)\in \Sj\times\Sj\).

Next, by Fubini's theorem and the previous normalization identity,
\[
\int_{\Sj}q_{j,h_j}(\ba,\bb)\,\rd \bb
=
\int_{\Sj}\int_{\Sj}\varphi_{j,h_j}(\ba,\bu)\varphi_{j,h_j}(\bb,\bu)\,\rd \bu\,\rd \bb
=
\int_{\Sj}\varphi_{j,h_j}(\ba,\bu)\,\rd \bu
=
1.
\]
Thus,
\[
\int_{\Sj}\int_{\Sj}q_{j,h_j}(\ba,\bb)^2\,\rd \ba\,\rd \bb
\leq
\lrp{\sup_{\ba,\bb}q_{j,h_j}(\ba,\bb)}
\int_{\Sj}\int_{\Sj}q_{j,h_j}(\ba,\bb)\,\rd \ba\,\rd \bb
=
O(h_j^{-d_j}).
\]

Using the product structure of \(Q_n\),
\[
\int_{\Sr}\int_{\Sr}Q_n(\bx,\by)^2\,\rd \bx\,\rd \by
=
\prod_{j=1}^r
\int_{\Sj}\int_{\Sj}q_{j,h_j}(\ba,\bb)^2\,\rd \ba\,\rd \bb
=
O\lrp{\prod_{j=1}^r h_j^{-d_j}}
=
O\lrp{\rho(\bh)^{-1}}.
\]
Consequently,
\[
A_{n,1}=O\lrp{\rho(\bh)^{-1}},
\]
and therefore
\[
\mathsf{E}\lrc{G_n^2(\bX_1,\bX_2)}
\leq 2A_{n,1}+O(1)
=
O\lrp{\rho(\bh)^{-1}}.
\]

This proves the required bound for \(\mathsf{E}\lrc{G_n^2(\bX_1,\bX_2)}\).

\section{Variance of $I_{n,2}$}

\begin{proof}
Recall that
\[
I_{n,2}
=
\frac{1}{n^2}\sum_{i=1}^n Z_{n,i},
\qquad
Z_{n,i}
:=
\int_{S^d} \overline L_h(\bx,\bX_i)^2 \, d\bx,
\]
where
\[
\overline L_h(\bx,\bX_i)
=
L_h(\bx,\bX_i)-E\{L_h(\bx,\bX_i)\}.
\]
Since \(Z_{n,1},\ldots,Z_{n,n}\) are iid,
\[
\Var(I_{n,2})
=
\Var\!\left(\frac{1}{n^2}\sum_{i=1}^n Z_{n,i}\right)
=
\frac{1}{n^4}\sum_{i=1}^n \Var(Z_{n,i})
=
\frac{1}{n^3}\Var(Z_{n,1})
\leq
\frac{1}{n^3}E(Z_{n,1}^2).
\]
Therefore, it is enough to prove that
\[
E(Z_{n,1}^2)=O\!\left(\rho(h)^{-2}\right).
\]

Write
\[
m_h(\bx):=E\{L_h(\bx,\bX_1)\}.
\]
Then
\[
Z_{n,1}
=
\int_{S^d}\bigl(L_h(\bx,\bX_1)-m_h(\bx)\bigr)^2\,d\bx
\leq
2\int_{S^d}L_h(\bx,\bX_1)^2\,d\bx
+
2\int_{S^d}m_h(\bx)^2\,d\bx.
\]
Hence, by \((a+b)^2\leq 2a^2+2b^2\),
\[
Z_{n,1}^2
\leq
8\left(\int_{S^d}L_h(\bx,\bX_1)^2\,d\bx\right)^2
+
8\left(\int_{S^d}m_h(\bx)^2\,d\bx\right)^2.
\]
Thus,
\[
E(Z_{n,1}^2)
\leq
8E\left[\left(\int_{S^d}L_h(\bx,\bX_1)^2\,d\bx\right)^2\right]
+
8\left(\int_{S^d}m_h(\bx)^2\,d\bx\right)^2.
\]

For the first term, note that, for any fixed \(\by\in S^d\),
\[
\int_{S^d}L_h(\bx,\by)^2\,d\bx
=
c_{d,L}(h)^2
\int_{S^d}
L\!\left(
\frac{1-\bx_1^\top\by_1}{h_1^2},
\ldots,
\frac{1-\bx_r^\top\by_r}{h_r^2}
\right)^2
\,d\bx.
\]
By the same change of variables used in the definition of \(c_{d,L}(h)\),
\[
\int_{S^d}
L\!\left(
\frac{1-\bx_1^\top\by_1}{h_1^2},
\ldots,
\frac{1-\bx_r^\top\by_r}{h_r^2}
\right)^2
\,d\bx
=
\rho(h)\lambda_{h,d}(L^2),
\]
and, since \(c_{d,L}(h)^{-1}=\rho(h)\lambda_{h,d}(L)\),
\[
\int_{S^d}L_h(\bx,\by)^2\,d\bx
=
\frac{\lambda_{h,d}(L^2)}{\rho(h)\lambda_{h,d}(L)^2}
=
O\!\left(\rho(h)^{-1}\right).
\]
Therefore,
\[
E\left[\left(\int_{S^d}L_h(\bx,\bX_1)^2\,d\bx\right)^2\right]
=
O\!\left(\rho(h)^{-2}\right).
\]

For the second term, since \(f\) is bounded on the compact set \(S^d\),
\[
m_h(\bx)
=
\int_{S^d}L_h(\bx,\by)f(\by)\,d\by
\leq
\|f\|_\infty \int_{S^d}L_h(\bx,\by)\,d\by
=
\|f\|_\infty.
\]
Hence,
\[
\int_{S^d}m_h(\bx)^2\,d\bx
\leq
\|f\|_\infty^2 \int_{S^d}d\bx
=
\omega_d\|f\|_\infty^2
=
O(1),
\]
so that
\[
\left(\int_{S^d}m_h(\bx)^2\,d\bx\right)^2
=
O(1).
\]

Combining the two bounds,
\[
E(Z_{n,1}^2)
=
O\!\left(\rho(h)^{-2}\right)+O(1)
=
O\!\left(\rho(h)^{-2}\right),
\]
because \(\rho(h)\to 0\). Consequently,
\[
\Var(I_{n,2})
\leq
\frac{1}{n^3}E(Z_{n,1}^2)
=
O\!\left(n^{-3}\rho(h)^{-2}\right).
\]
This concludes the proof.
\end{proof}

\section{Plots}

\begin{figure}
    \centering
    \includegraphics[width=1\linewidth]{testeplot.pdf}
    \caption{Caption}
    \label{fig:placeholder}
\end{figure}

\section{New definitions}

\[
\lambda_{d_j}(L)
:=
2^{d_j/2-1}\omega_{d_j-1}
\int_0^\infty L(s)s^{d_j/2-1}\,ds,
\qquad
\lambda_{\bd}(L)
:=
\prod_{j=1}^r \lambda_{d_j}(L).
\]

\[
v_{\bd}(L)
:=
\frac{\lambda_{\bd}(L^2)}{\lambda_{\bd}(L)^2}
=
\prod_{j=1}^r
\frac{\lambda_{d_j}(L^2)}{\lambda_{d_j}(L)^2}.
\]

\bibliographystyle{alpha}
\bibliography{sample}

\end{document}

O()